% =========================================================
% Chapter: Ray Optics, Reflection, Refraction, Mirrors, Lenses
% POSN 1 Physics Preparation Notes
% Language: Thai
% Compiler: XeLaTeX
%
% Required environments/packages from the main project:
%   graphicx, siunitx, tcolorbox
%   examplebox, exercisebox
% =========================================================

\chapter{แสง การสะท้อน การหักเห และการเกิดภาพด้วยกระจกและเลนส์}

\begin{center}
    {\Large \textbf{Ray Optics and Image Formation}}\\[4pt]
    {\normalsize เอกสารเตรียมสอบคัดเลือก สอวน. ฟิสิกส์ ค่าย 1}\\[6pt]
    {\normalsize จัดทำโดย \textbf{Tames Thanakrit Damduan}}\\[2pt]
    {\footnotesize \textcopyright\ 2026 Tames Thanakrit Damduan. All rights reserved.}\\[8pt]
\end{center}

\sloppy

\section*{เป้าหมายของบทเรียน}

หลังจากเรียนบทนี้ นักเรียนควรทำสิ่งต่อไปนี้ได้

\begin{enumerate}
    \item อธิบายแบบจำลองรังสีของแสงและเงื่อนไขที่ใช้ทัศนศาสตร์เชิงรังสีได้
    \item ใช้กฎการสะท้อนและกฎของสเนลล์หาทิศทางของแสงได้
    \item เชื่อมโยงดรรชนีหักเหกับอัตราเร็วและความยาวคลื่นของแสงได้
    \item วิเคราะห์มุมวิกฤต การสะท้อนกลับหมด และมุมบริวสเตอร์ได้
    \item วิเคราะห์การกระจายแสงด้วยปริซึมและความแตกต่างของดรรชนีหักเหตามสีได้
    \item ใช้ข้อตกลงเครื่องหมายสำหรับกระจกและเลนส์ได้อย่างสม่ำเสมอ
    \item อนุมานสมการกระจกทรงกลมและความสัมพันธ์ $f=R/2$ ได้
    \item อนุมานสมการการหักเหที่ผิวทรงกลมและสมการช่างทำเลนส์ได้
    \item ใช้สมการกระจก สมการเลนส์ และกำลังขยายเพื่อหาตำแหน่งและลักษณะภาพได้
    \item วิเคราะห์ระบบหลายชิ้น เช่น เลนส์สองอัน ระบบเลนส์กับกระจก และระบบขยายลำแสงได้
    \item ใช้การเปลี่ยนแปลงเล็กน้อยวิเคราะห์ผลของรัศมีความโค้งและดรรชนีหักเหต่อความยาวโฟกัสได้
    \item แก้ข้อสอบจริง สอวน. ที่เกี่ยวข้องกับทัศนศาสตร์เชิงรังสีตั้งแต่ปี 2560 ถึง 2568 ได้
\end{enumerate}

\begin{tcolorbox}[title={แนวคิดสำคัญของบทนี้}]
แสงมีสมบัติเป็นคลื่นแม่เหล็กไฟฟ้า แต่เมื่อความยาวคลื่นเล็กมากเมื่อเทียบกับขนาดของช่องเปิด ผิวสะท้อน และเลนส์ เราสามารถแทนการเดินทางของแสงด้วยเส้นตรงที่เรียกว่า \textbf{รังสีแสง}

โจทย์ส่วนใหญ่ในบทนี้จึงลดรูปเป็นสามขั้น

\begin{enumerate}
    \item หาทิศของรังสีด้วยกฎการสะท้อนหรือกฎของสเนลล์
    \item หาตำแหน่งภาพด้วยสมการกระจกหรือสมการเลนส์
    \item ตรวจลักษณะภาพด้วยเครื่องหมายของระยะภาพและกำลังขยาย
\end{enumerate}
\end{tcolorbox}

% =========================================================
% =========================================================
\section{แบบจำลองรังสีของแสงและที่มาของกฎพื้นฐาน}

\begin{tcolorbox}[title={ก่อนเริ่ม: สมการในบทนี้มีสามชนิด}]
ในบทนี้จะไม่เขียนสูตรใหม่แล้วปล่อยให้นักเรียนท่องโดยไม่ทราบที่มา เราจะแยกสมการเป็นสามชนิดอย่างชัดเจน


\end{tcolorbox}

\subsection{เมื่อใดจึงใช้ทัศนศาสตร์เชิงรังสีได้}

แสงเป็นคลื่นแม่เหล็กไฟฟ้า แต่ถ้าความยาวคลื่น $\lambda$ เล็กมากเมื่อเทียบกับขนาดลักษณะเฉพาะของระบบ $L$ เช่น ขนาดกระจก เส้นผ่านศูนย์กลางเลนส์ หรือความกว้างของช่องเปิด

\[
\lambda\ll L
\]

การเลี้ยวเบนมักมีผลน้อย เราจึงประมาณเส้นทางการส่งผ่านพลังงานด้วยเส้นที่เรียกว่า \textbf{รังสีแสง}

\begin{tcolorbox}[title={แบบจำลองรังสี}]
รังสีแสงคือเส้นที่ชี้ทิศการแพร่ของพลังงานแสง ในตัวกลางเนื้อเดียวรังสีตั้งฉากกับหน้าคลื่นและเป็นเส้นตรง
\end{tcolorbox}

เงื่อนไข $\lambda\ll L$ เป็นเงื่อนไขของการประมาณ ไม่ใช่กฎที่บอกว่าแสงไม่มีสมบัติคลื่น เมื่อ $\lambda$ มีขนาดเทียบได้กับ $L$ ต้องกลับไปใช้ทัศนศาสตร์เชิงคลื่น

\subsection{ดรรชนีหักเหและความยาวคลื่นในตัวกลาง}

ในสุญญากาศ แสงมีอัตราเร็ว

\[
c\approx 3.00\times10^8\,\mathrm{m\,s^{-1}}
\]

ถ้าแสงเดินทางในตัวกลางด้วยอัตราเร็ว $v$ เรา \textbf{นิยาม} ดรรชนีหักเหเป็น

\begin{tcolorbox}[title={นิยาม: ดรรชนีหักเห}]
\[
\boxed{n\equiv\frac{c}{v}}
\]
\end{tcolorbox}

จากนิยาม คูณทั้งสองข้างด้วย $v$ แล้วหารด้วย $n$

\[
nv=c
\]

\[
\boxed{v=\frac{c}{n}}
\]

ดังนั้น $n$ มากหมายถึงแสงเดินทางช้าลงในตัวกลางนั้น

ต่อไปจะหาความยาวคลื่นในตัวกลาง ไม่ใช่ประกาศสูตรขึ้นมาใหม่ เริ่มจากความสัมพันธ์ของคลื่น

\[
v=f\lambda
\]

ความถี่ของคลื่นถูกกำหนดโดยแหล่งกำเนิด และที่รอยต่อหน้าคลื่นต้องต่อเนื่องกัน จึงมีความถี่เดียวกันทั้งสองด้าน

\[
f_{\text{medium}}=f_{\text{vacuum}}=f
\]

ในสุญญากาศ

\[
c=f\lambda_0
\]

จึงได้

\[
\lambda_0=\frac{c}{f}
\]

ในตัวกลาง

\[
\lambda=\frac{v}{f}
\]

แทน $v=c/n$

\[
\lambda=\frac{c/n}{f}
=\frac{1}{n}\frac{c}{f}
\]

แต่ $c/f=\lambda_0$ จึงได้

\begin{tcolorbox}[title={ความยาวคลื่นในตัวกลาง}]
\[
\boxed{\lambda=\frac{\lambda_0}{n}}
\]
\end{tcolorbox}

ข้อสรุปสำคัญคือ เมื่อแสงเปลี่ยนตัวกลาง ความถี่คงเดิม แต่อัตราเร็วและความยาวคลื่นเปลี่ยน

\subsection{เวลาเดินทางและระยะทางเชิงแสง}

ถ้าแสงเดินทางระยะ $L$ ในตัวกลางที่มีอัตราเร็ว $v$ เวลาที่ใช้คือ

\[
T=\frac{L}{v}
\]

แทน $v=c/n$

\[
T=\frac{L}{c/n}=\frac{nL}{c}
\]

ปริมาณ $nL$ เรียกว่า \textbf{ระยะทางเชิงแสง} สำหรับหลายช่วง

\[
T=\frac{1}{c}\sum_i n_iL_i
\]

ในทัศนศาสตร์เชิงรังสี เส้นทางจริงเป็นเส้นทางที่เวลาเดินทางมีค่าอยู่กับที่ต่อการเปลี่ยนเส้นทางเล็กน้อย ในกรณีง่ายของบทนี้ค่านี้เป็นค่าน้อยที่สุด หลักนี้เป็นผลระดับรังสีของการรวมเฟสของคลื่น และจะใช้เป็นจุดตั้งต้นเพื่ออนุมานกฎการสะท้อนและการหักเห

% =========================================================
\section{การสะท้อนของแสง}

\subsection{อนุมานกฎการสะท้อนจากเวลาเดินทาง}

ให้กระจกเงาราบอยู่ตามแกน $x$ และให้แสงเดินจากจุด $A$ ไปสะท้อนที่จุด $P$ แล้วไปยังจุด $B$ จุดทั้งสองอยู่ด้านเดียวกับกระจก กำหนดพิกัด

\[
A=(-a,y_A),\qquad P=(x,0),\qquad B=(b,y_B)
\]

ระยะจาก $A$ ถึง $P$ คือ

\[
AP=\sqrt{(x+a)^2+y_A^2}
\]

ระยะจาก $P$ ถึง $B$ คือ

\[
PB=\sqrt{(b-x)^2+y_B^2}
\]

แสงเดินทางอยู่ในตัวกลางเดียวกัน อัตราเร็วคงตัว ดังนั้นทำให้เวลาอยู่กับที่เท่ากับทำให้ระยะรวม

\[
L(x)=AP+PB
\]

อยู่กับที่ หาอนุพันธ์ตามตำแหน่งสะท้อน $x$

\[
\frac{dL}{dx}
=
\frac{x+a}{AP}
-
\frac{b-x}{PB}
\]

เส้นทางจริงต้องมี

\[
\frac{dL}{dx}=0
\]

จึงได้

\[
\frac{x+a}{AP}
=
\frac{b-x}{PB}
\]

จากสามเหลี่ยมด้านซ้าย

\[
\frac{x+a}{AP}=\sin\theta_i
\]

และจากสามเหลี่ยมด้านขวา

\[
\frac{b-x}{PB}=\sin\theta_r
\]

ดังนั้น

\[
\sin\theta_i=\sin\theta_r
\]

เมื่อ $0\le\theta_i,\theta_r<90^\circ$ ฟังก์ชันไซน์เพิ่มขึ้นตลอดช่วง จึงสรุปได้ว่า

\begin{tcolorbox}[title={กฎการสะท้อน}]
\[
\boxed{\theta_i=\theta_r}
\]

มุมทั้งสองต้องวัดจากเส้นปกติต่อผิว
\end{tcolorbox}

รังสีตกกระทบ รังสีสะท้อน และเส้นปกติอยู่ในระนาบเดียวกัน เพราะการเบี่ยงออกนอกระนาบที่เชื่อม $A$, $B$ และเส้นปกติจะเพิ่มระยะทางโดยไม่ช่วยลดระยะส่วนอื่น

\subsection{การสะท้อนแบบเป็นระเบียบและแบบกระจาย}

ผิวเรียบมีเส้นปกติใกล้เคียงกันทุกจุด รังสีขนานจึงสะท้อนออกไปอย่างเป็นระเบียบ

ผิวขรุขระมีเส้นปกติแต่ละจุดต่างกัน รังสีจึงออกหลายทิศ อย่างไรก็ตามทุกจุดยังคงใช้

\[
\theta_i=\theta_r
\]

ความแตกต่างจึงไม่ได้เกิดจากกฎเปลี่ยน แต่เกิดจากแนวเส้นปกติของผิวเปลี่ยนไปในแต่ละตำแหน่ง

\subsection{ภาพจากกระจกเงาราบ}

ให้กระจกอยู่ที่ $x=0$ และวัตถุจุดอยู่ที่

\[
O=(-p,h)
\]

สะท้อนพิกัดของจุด $O$ ข้ามระนาบกระจก จะได้จุด

\[
I=(+p,h)
\]

จากกฎการสะท้อน เส้นทางหัก $O\to P\to$ ผู้สังเกต มีทิศเดียวกับเส้นตรงที่ต่อย้อนหลังไปยัง $I$ สำหรับทุกจุดสะท้อน $P$ ผู้สังเกตจึงเห็นรังสีเสมือนออกมาจาก $I$

ระยะเชิงขนาดของภาพเท่ากับระยะวัตถุ

\[
|q|=p
\]

แต่ภาพอยู่หลังกระจกและเป็นภาพเสมือน ตามข้อตกลงภาพเสมือนต้องมี $q<0$ จึงได้

\[
\boxed{q=-p}
\]

พิกัดแนวดิ่งไม่เปลี่ยนจากการสะท้อนข้ามกระจก

\[
h'=h
\]

จากนิยามกำลังขยาย $m\equiv h'/h$

\[
m=\frac{h}{h}=1
\]

ดังนั้น

\begin{tcolorbox}[title={กระจกเงาราบ}]
\[
\boxed{q=-p}
\qquad
\boxed{m=+1}
\]
\end{tcolorbox}

ภาพจึงเสมือน ตั้งตรง และมีขนาดเท่าวัตถุ การที่ดูเหมือนสลับซ้ายขวาเกิดจากการสะท้อนพิกัดในทิศตั้งฉากกับกระจก ไม่ใช่การหมุนภาพรอบแกนตั้ง

% =========================================================
\section{การหักเหของแสง}

\subsection{อนุมานกฎของสเนลล์จากเวลาเดินทาง}

ให้รอยต่อระหว่างตัวกลางอยู่ตามแกน $x$ ตัวกลาง 1 อยู่ด้านบนมีอัตราเร็ว $v_1$ และตัวกลาง 2 อยู่ด้านล่างมีอัตราเร็ว $v_2$ แสงเดินจาก

\[
A=(-a,h_1)
\]

ผ่านรอยต่อที่

\[
P=(x,0)
\]

ไปยัง

\[
B=(b,-h_2)
\]

เวลารวมเป็น

\[
T(x)
=
\frac{\sqrt{(x+a)^2+h_1^2}}{v_1}
+
\frac{\sqrt{(b-x)^2+h_2^2}}{v_2}
\]

หาอนุพันธ์

\[
\frac{dT}{dx}
=
\frac{x+a}{v_1AP}
-
\frac{b-x}{v_2PB}
\]

เส้นทางจริงมี $dT/dx=0$

\[
\frac{x+a}{AP}\frac{1}{v_1}
=
\frac{b-x}{PB}\frac{1}{v_2}
\]

จากเรขาคณิต

\[
\frac{x+a}{AP}=\sin\theta_1
\qquad
\frac{b-x}{PB}=\sin\theta_2
\]

จึงได้

\[
\frac{\sin\theta_1}{v_1}
=
\frac{\sin\theta_2}{v_2}
\]

ใช้ $v_i=c/n_i$

\[
\frac{\sin\theta_1}{c/n_1}
=
\frac{\sin\theta_2}{c/n_2}
\]

\[
\frac{n_1\sin\theta_1}{c}
=
\frac{n_2\sin\theta_2}{c}
\]

คูณด้วย $c$

\begin{tcolorbox}[title={กฎของสเนลล์}]
\[
\boxed{n_1\sin\theta_1=n_2\sin\theta_2}
\]
\end{tcolorbox}

ถ้า $n_2>n_1$

\[
\sin\theta_2=\frac{n_1}{n_2}\sin\theta_1<\sin\theta_1
\]

สำหรับมุมระหว่าง $0$ ถึง $90^\circ$ จึงมี

\[
\theta_2<\theta_1
\]

แสงหักเหเข้าใกล้เส้นปกติ ในทางกลับกัน ถ้า $n_2<n_1$ แสงหักเหออกจากเส้นปกติ

\subsection{ความลึกปรากฏ}

วัตถุอยู่ในตัวกลาง 1 ลึกจริง $d$ ผู้สังเกตอยู่ในตัวกลาง 2 และมองผ่านผิวราบ ให้รังสีหนึ่งตัดผิวห่างจากแนวตั้งเป็นระยะ $x$

จากสามเหลี่ยมในตัวกลาง 1

\[
\tan\theta_1=\frac{x}{d}
\]

เมื่อต่อรังสีหักเหย้อนหลังไปยังภาพเสมือนที่ลึก $d_{\text{app}}$

\[
\tan\theta_2=\frac{x}{d_{\text{app}}}
\]

สมการสองบรรทัดนี้เป็นเรขาคณิตที่แน่นอน แต่กฎของสเนลล์เกี่ยวข้องกับไซน์ ไม่ใช่แทนเจนต์ ดังนั้นความลึกปรากฏคงที่แบบง่ายเกิดขึ้นภายใต้การมองเกือบตั้งฉากเท่านั้น

สำหรับมุมเล็กในหน่วยเรเดียน

\[
\sin\theta\approx\tan\theta\approx\theta
\]

จึงมี

\[
\sin\theta_1\approx\frac{x}{d}
\qquad
\sin\theta_2\approx\frac{x}{d_{\text{app}}}
\]

แทนในกฎของสเนลล์

\[
n_1\frac{x}{d}
=
n_2\frac{x}{d_{\text{app}}}
\]

ตัด $x\ne0$

\[
\frac{n_1}{d}
=
\frac{n_2}{d_{\text{app}}}
\]

คูณไขว้

\begin{tcolorbox}[title={ความลึกปรากฏเมื่อมองเกือบตั้งฉาก}]
\[
\boxed{d_{\text{app}}=\frac{n_2}{n_1}d}
\]
\end{tcolorbox}

เมื่อน้ำเป็นตัวกลาง 1 และอากาศเป็นตัวกลาง 2 จะมี $n_2<n_1$ จึงเห็นวัตถุตื้นกว่าความจริง

\subsection{การสะท้อนกลับหมดและมุมวิกฤต}

จากกฎของสเนลล์

\[
\sin\theta_2=\frac{n_1}{n_2}\sin\theta_1
\]

ถ้า $n_1\le n_2$ ด้านขวาไม่เกิน 1 สำหรับทุกมุมตกกระทบ จึงมีรังสีหักเหได้เสมอ

ถ้า $n_1>n_2$ เมื่อเพิ่ม $\theta_1$ ค่า $\sin\theta_2$ จะเพิ่มจนถึงค่าสูงสุด 1 ขอบเขตนี้เรียกว่ามุมวิกฤต $\theta_c$ และรังสีหักเหเลียบผิว

\[
\theta_2=90^\circ
\]

แทนในกฎของสเนลล์

\[
n_1\sin\theta_c=n_2\sin90^\circ
\]

เพราะ $\sin90^\circ=1$

\begin{tcolorbox}[title={มุมวิกฤต}]
\[
\boxed{\sin\theta_c=\frac{n_2}{n_1}}
\qquad (n_1>n_2)
\]
\end{tcolorbox}

ถ้า $\theta_i>\theta_c$ สมการจะต้องการ $\sin\theta_2>1$ ซึ่งเป็นไปไม่ได้สำหรับมุมจริง จึงไม่มีคลื่นเดินทางออกไปในตัวกลาง 2 และเกิดการสะท้อนกลับหมด

\subsection{มุมบริวสเตอร์}

มุมบริวสเตอร์ไม่ใช่ผลจากกฎของสเนลล์เพียงสมการเดียว ต้องใช้ผลของคลื่นแม่เหล็กไฟฟ้าที่รอยต่อด้วย สำหรับแสงโพลาไรซ์ในระนาบตกกระทบ สัมประสิทธิ์แอมพลิจูดการสะท้อนที่ได้จากเงื่อนไขขอบเขตของสนามไฟฟ้าและสนามแม่เหล็กคือ

\[
r_p=
\frac{n_2\cos\theta_i-n_1\cos\theta_2}
{n_2\cos\theta_i+n_1\cos\theta_2}
\]

ที่มุมบริวสเตอร์ แสงโพลาไรซ์ชนิดนี้ไม่สะท้อน จึงมี

\[
r_p=0
\]

เศษของเศษส่วนต้องเป็นศูนย์

\[
n_2\cos\theta_B=n_1\cos\theta_2
\]

จากกฎของสเนลล์

\[
n_1\sin\theta_B=n_2\sin\theta_2
\]

เขียนอัตราส่วน $n_2/n_1$ จากทั้งสองสมการ

\[
\frac{n_2}{n_1}
=
\frac{\cos\theta_2}{\cos\theta_B}
\]

และ

\[
\frac{n_2}{n_1}
=
\frac{\sin\theta_B}{\sin\theta_2}
\]

เท่ากันจึงได้

\[
\sin\theta_B\cos\theta_B
=
\sin\theta_2\cos\theta_2
\]

หรือ

\[
\sin(2\theta_B)=\sin(2\theta_2)
\]

สำหรับรอยต่อที่ $n_1\ne n_2$ คำตอบที่เกี่ยวข้องคือ

\[
\theta_B+\theta_2=90^\circ
\]

ดังนั้น $\sin\theta_2=\cos\theta_B$ แทนในกฎของสเนลล์

\[
n_1\sin\theta_B=n_2\cos\theta_B
\]

หารด้วย $n_1\cos\theta_B$

\begin{tcolorbox}[title={กฎของบริวสเตอร์}]
\[
\boxed{\tan\theta_B=\frac{n_2}{n_1}}
\]
\end{tcolorbox}

\subsection{การกระจายแสงและปริซึมบาง}

ดรรชนีหักเหขึ้นกับความยาวคลื่น โดยทั่วไปในแก้วช่วงแสงที่ตามองเห็น

\[
n_{\text{blue}}>n_{\text{red}}
\]

จึงทำให้สีน้ำเงินเบี่ยงเบนมากกว่าสีแดง

ให้ปริซึมมีมุมยอด $A$ มุมตกกระทบและหักเหที่ผิวแรกเป็น $i_1,r_1$ และที่ผิวที่สองเป็น $r_2,i_2$ ตามลำดับ

จากเรขาคณิตของปริซึม

\[
r_1+r_2=A
\]

การเปลี่ยนทิศที่ผิวแรกเท่ากับ

\[
\delta_1=i_1-r_1
\]

การเปลี่ยนทิศที่ผิวที่สองเท่ากับ

\[
\delta_2=i_2-r_2
\]

ดังนั้นมุมเบี่ยงเบนรวม

\[
\delta=\delta_1+\delta_2
\]

\[
\delta=i_1+i_2-(r_1+r_2)
\]

ใช้ $r_1+r_2=A$

\[
\delta=i_1+i_2-A
\]

สำหรับปริซึมมุมเล็ก ทุกมุมมีค่าน้อยและใช้ $\sin\theta\approx\theta$ ได้ กฎของสเนลล์ที่ผิวแรกให้

\[
\sin i_1=n\sin r_1
\]

จึงประมาณเป็น

\[
i_1\approx nr_1
\]

ที่ผิวที่สอง

\[
n\sin r_2=\sin i_2
\]

จึงประมาณเป็น

\[
i_2\approx nr_2
\]

แทนในสมการมุมเบี่ยงเบน

\[
\delta\approx nr_1+nr_2-A
\]

\[
\delta\approx n(r_1+r_2)-A
\]

ใช้ $r_1+r_2=A$

\begin{tcolorbox}[title={มุมเบี่ยงเบนของปริซึมบาง}]
\[
\boxed{\delta\approx(n-1)A}
\]
\end{tcolorbox}

สำหรับสองสีที่มีดรรชนี $n_a$ และ $n_b$

\[
\delta_a\approx(n_a-1)A
\]

\[
\delta_b\approx(n_b-1)A
\]

ลบกัน

\begin{tcolorbox}[title={มุมแยกของสองสีในปริซึมบาง}]
\[
\boxed{\Delta\delta\approx(n_a-n_b)A}
\]
\end{tcolorbox}

ถ้ามุมไม่เล็กต้องกลับไปใช้กฎของสเนลล์แบบเต็ม ไม่ควรใช้ $\delta\approx(n-1)A$

% =========================================================
\section{ข้อตกลงเครื่องหมายสำหรับการเกิดภาพ}

บทนี้ใช้ข้อตกลงแบบ \textbf{วัตถุจริงและภาพจริงเป็นบวก} เพื่อให้สมการกระจกและเลนส์มีรูปเดียวกัน

\begin{tcolorbox}[title={ข้อตกลงเครื่องหมายที่ใช้ในบทนี้}]
\textbf{ระยะวัตถุ $p$}

\[
p>0
\]

เมื่อรังสีที่เข้าสู่อุปกรณ์แยกออกมาจากวัตถุจริง

\[
p<0
\]

เมื่อรังสีที่เข้าสู่อุปกรณ์กำลังรวมเข้าหาจุดซึ่งทำหน้าที่เป็นวัตถุเสมือน

\medskip
\textbf{ระยะภาพ $q$}

\[
q>0
\]

เมื่อรังสีหลังออกจากอุปกรณ์ตัดกันจริง

\[
q<0
\]

เมื่อรังสีหลังออกจากอุปกรณ์แยกจากกัน แต่เส้นต่อย้อนหลังตัดกัน

\medskip
\textbf{ความยาวโฟกัส $f$}

\[
f>0
\]

สำหรับอุปกรณ์รวมแสง เช่น กระจกเว้าและเลนส์นูนในอากาศ

\[
f<0
\]

สำหรับอุปกรณ์กระจายแสง เช่น กระจกนูนและเลนส์เว้าในอากาศ
\end{tcolorbox}

กำลังขยายตามขวางเริ่มจาก \textbf{นิยาม}

\begin{tcolorbox}[title={นิยาม: กำลังขยายตามขวาง}]
\[
\boxed{m\equiv\frac{h'}{h}}
\]
\end{tcolorbox}

เครื่องหมายของ $m$ แปลผลดังนี้

\[
m>0 \quad\Rightarrow\quad \text{ภาพตั้งตรง}
\]

\[
m<0 \quad\Rightarrow\quad \text{ภาพกลับหัว}
\]

ส่วนความสัมพันธ์ $m=-q/p$ ยังไม่ถือเป็นนิยาม เราจะอนุมานจากสามเหลี่ยมคล้ายแยกสำหรับกระจกและเลนส์

% =========================================================
\section{กระจกทรงกลม}

\subsection{ส่วนประกอบและเงื่อนไขรังสีใกล้แกน}

ให้ $P$ เป็นขั้วกระจก $C$ เป็นศูนย์กลางความโค้ง และ

\[
R=PC
\]

เป็นรัศมีความโค้ง การอนุมานต่อไปใช้รังสีที่กระทบใกล้แกนมุขสำคัญ ทำให้มุมทุกมุมมีค่าน้อยและใช้

\[
\sin\theta\approx\tan\theta\approx\theta
\]

ได้

\subsection{อนุมานสมการกระจกทรงกลม}

พิจารณากระจกเว้า วัตถุอยู่บนแกนที่ระยะ $p$ และภาพอยู่บนแกนที่ระยะ $q$ รังสีหนึ่งตกกระจกที่จุด $Q$ ซึ่งสูงจากแกนเป็น $h$

กำหนดมุมเล็ก

\[
\alpha=\text{มุมระหว่างรังสีตกกระทบกับแกน}
\]

\[
\beta=\text{มุมระหว่างรังสีสะท้อนกับแกน}
\]

\[
\gamma=\text{มุมระหว่างเส้นปกติ }QC\text{ กับแกน}
\]

จากรูปเรขาคณิตของกรณีวัตถุอยู่นอกศูนย์กลางความโค้ง มุมตกกระทบและมุมสะท้อนเขียนได้เป็น

\[
\theta_i=\gamma-\alpha
\]

\[
\theta_r=\beta-\gamma
\]

ใช้กฎการสะท้อน $\theta_i=\theta_r$

\[
\gamma-\alpha=\beta-\gamma
\]

ย้ายข้าง

\[
\alpha+\beta=2\gamma
\]

สำหรับมุมเล็ก จากสามเหลี่ยมที่มีความสูงร่วม $h$

\[
\tan\alpha=\frac{h}{p}
\quad\Rightarrow\quad
\alpha\approx\frac{h}{p}
\]

\[
\tan\beta=\frac{h}{q}
\quad\Rightarrow\quad
\beta\approx\frac{h}{q}
\]

\[
\tan\gamma=\frac{h}{R}
\quad\Rightarrow\quad
\gamma\approx\frac{h}{R}
\]

แทนใน $\alpha+\beta=2\gamma$

\[
\frac{h}{p}+\frac{h}{q}=2\frac{h}{R}
\]

ตัด $h$

\[
\boxed{\frac{1}{p}+\frac{1}{q}=\frac{2}{R}}
\]

ผลนี้มาจากกฎการสะท้อนและการประมาณรังสีใกล้แกน ยังไม่ได้ใช้ความยาวโฟกัส

\subsection{อนุมานความสัมพันธ์ $f=R/2$}

จุดโฟกัสภาพคือจุดที่รังสีขนานแกนรวมกัน รังสีขนานเทียบเท่ากับวัตถุที่อยู่ไกลมาก

\[
p\to\infty
\]

ดังนั้น

\[
\frac{1}{p}\to0
\]

และภาพเกิดที่ $q=f$ แทนในสมการที่เพิ่งอนุมาน

\[
0+\frac{1}{f}=\frac{2}{R}
\]

กลับเศษส่วน

\[
\boxed{f=\frac{R}{2}}
\]

จึงเขียน $2/R=1/f$ และได้

\begin{tcolorbox}[title={สมการกระจกทรงกลม}]
\[
\boxed{\frac{1}{p}+\frac{1}{q}=\frac{1}{f}}
\qquad
\boxed{f=\frac{R}{2}}
\]
\end{tcolorbox}

สังเกตว่า $f=R/2$ ไม่ใช่นิยามลอย ๆ แต่เป็นผลจากสมการกระจกเมื่อรังสีตกกระทบขนาน

\subsection{อนุมานกำลังขยายของกระจก}

ให้รังสีจากยอดวัตถุตกที่ขั้วกระจก $P$ เส้นปกติที่ขั้วคือแกนมุขสำคัญ กฎการสะท้อนทำให้มุมของรังสีก่อนและหลังสะท้อนกับแกนมีขนาดเท่ากัน

จากสามเหลี่ยมด้านวัตถุ

\[
\tan\theta=\frac{h}{p}
\]

จากสามเหลี่ยมด้านภาพ ภาพจริงกลับหัวจึงมีความสูงเชิงเครื่องหมาย $h'<0$

\[
\tan\theta=-\frac{h'}{q}
\]

เท่ากันจึงได้

\[
\frac{h}{p}=-\frac{h'}{q}
\]

จัดรูป

\[
\frac{h'}{h}=-\frac{q}{p}
\]

ใช้คำนิยาม $m=h'/h$

\begin{tcolorbox}[title={กำลังขยายของกระจกทรงกลม}]
\[
\boxed{m=-\frac{q}{p}}
\]
\end{tcolorbox}

\subsection{รังสีสำคัญของกระจกมาจากสมการใด}

\begin{enumerate}
    \item รังสีขนานแกนมี $p\to\infty$ จึงได้ $q=f$ และสะท้อนผ่านจุดโฟกัส
    \item เพราะเส้นทางแสงย้อนกลับได้ รังสีที่มุ่งผ่านจุดโฟกัสจึงสะท้อนขนานแกน
    \item รังสีที่ผ่าน $C$ เดินตามแนวรัศมี จึงตกกระทบตามเส้นปกติ มี $\theta_i=0$ และสะท้อนย้อนแนวเดิม
    \item รังสีตกที่ขั้วกระจกใช้เส้นปกติเป็นแกนและใช้ $\theta_i=\theta_r$
\end{enumerate}

\subsection{กรณีภาพของกระจกเว้าและกระจกนูน}

ตารางต่อไปไม่ได้เป็นสูตรแยกชุดใหม่ แต่ได้จากเครื่องหมายของคำตอบ $q$ และ $m$ ในสมการกระจก

\begin{center}
\begin{tabular}{|c|c|c|c|}
\hline
ตำแหน่งวัตถุ & ตำแหน่งภาพ & ลักษณะภาพ & ขนาด \\
\hline
$p>2f$ & $f<q<2f$ & จริง กลับหัว & เล็กลง \\
\hline
$p=2f$ & $q=2f$ & จริง กลับหัว & เท่าเดิม \\
\hline
$f<p<2f$ & $q>2f$ & จริง กลับหัว & ใหญ่ขึ้น \\
\hline
$p=f$ & $q\to\infty$ & รังสีออกขนาน & ไม่กำหนด \\
\hline
$0<p<f$ & $q<0$ & เสมือน ตั้งตรง & ใหญ่ขึ้น \\
\hline
\end{tabular}
\end{center}

สำหรับกระจกนูน $f<0$ และวัตถุจริง $p>0$

\[
\frac{1}{q}=\frac{1}{f}-\frac{1}{p}
\]

ทั้ง $1/f$ และ $-1/p$ เป็นลบ จึงมี

\[
\frac{1}{q}<0
\quad\Rightarrow\quad
q<0
\]

จาก $m=-q/p$ เมื่อ $q<0$ และ $p>0$ จะได้ $m>0$ ภาพจึงตั้งตรง และเพราะ $|q|<p$ จึงมี $|m|<1$

\begin{examplebox}{ตัวอย่าง: กระจกเว้า}
วัตถุอยู่หน้ากระจกเว้าความยาวโฟกัส $f$ ที่ระยะ $p=3f$

เริ่มจากสมการกระจก

\[
\frac{1}{p}+\frac{1}{q}=\frac{1}{f}
\]

แทน $p=3f$

\[
\frac{1}{3f}+\frac{1}{q}=\frac{1}{f}
\]

ย้าย $1/(3f)$

\[
\frac{1}{q}
=
\frac{1}{f}-\frac{1}{3f}
=
\frac{2}{3f}
\]

กลับเศษส่วน

\[
q=\frac{3f}{2}
\]

กำลังขยาย

\[
m=-\frac{q}{p}
=-\frac{3f/2}{3f}
=-\frac{1}{2}
\]

$q>0$ บอกว่าเป็นภาพจริง $m<0$ บอกว่ากลับหัว และ $|m|=1/2$ บอกว่าภาพสูงครึ่งหนึ่งของวัตถุ
\end{examplebox}

% =========================================================
\section{การหักเหที่ผิวทรงกลม}

\subsection{ข้อตกลงเครื่องหมายของรัศมีผิว}

สำหรับผิวหักเห ให้ทิศแสงเดินจากซ้ายไปขวา

\[
R>0
\]

เมื่อศูนย์กลางความโค้งอยู่ด้านขวาของผิว หรือด้านที่แสงหักเหออกไป และ

\[
R<0
\]

เมื่อศูนย์กลางความโค้งอยู่ด้านซ้าย หรือด้านที่แสงตกกระทบ

\subsection{อนุมานสมการผิวหักเหทรงกลม}

พิจารณารังสีจากวัตถุจริงในตัวกลางดรรชนี $n_1$ ผ่านผิวทรงกลมไปยังตัวกลางดรรชนี $n_2$ ผิวมีศูนย์กลางความโค้งด้านขวา จึงมี $R>0$ รังสีตกผิวที่ความสูง $h$

กำหนดมุมเล็ก

\[
\alpha\approx\frac{h}{p}
\]

เป็นมุมของรังสีตกกระทบกับแกน

\[
\beta\approx\frac{h}{q}
\]

เป็นขนาดมุมของรังสีหักเหกับแกน และ

\[
\gamma\approx\frac{h}{R}
\]

เป็นมุมของเส้นปกติซึ่งชี้ไปยังศูนย์กลางความโค้ง

จากรูปเรขาคณิต รังสีตกกระทบอยู่เหนือแกน ส่วนเส้นปกติเอียงลงหาแกน จึงมี

\[
\theta_1=\alpha+\gamma
\]

รังสีหักเหและเส้นปกติอยู่ด้านเดียวกันของแกน โดยในกรณีรังสีใกล้แกน

\[
\theta_2=\gamma-\beta
\]

ใช้กฎของสเนลล์

\[
n_1\sin\theta_1=n_2\sin\theta_2
\]

สำหรับมุมเล็ก $\sin\theta\approx\theta$

\[
n_1(\alpha+\gamma)
=
n_2(\gamma-\beta)
\]

กระจายวงเล็บ

\[
n_1\alpha+n_1\gamma
=
n_2\gamma-n_2\beta
\]

ย้าย $n_2\beta$ มาซ้าย และย้าย $n_1\gamma$ ไปขวา

\[
n_1\alpha+n_2\beta
=(n_2-n_1)\gamma
\]

แทน $\alpha\approx h/p$, $\beta\approx h/q$, $\gamma\approx h/R$

\[
n_1\frac{h}{p}
+n_2\frac{h}{q}
=(n_2-n_1)\frac{h}{R}
\]

ตัด $h$

\begin{tcolorbox}[title={การหักเหที่ผิวทรงกลม}]
\[
\boxed{
\frac{n_1}{p}
+
\frac{n_2}{q}
=
\frac{n_2-n_1}{R}}
\]
\end{tcolorbox}

เราอนุมานโดยวาดกรณี $p,q,R>0$ เพื่อให้เห็นมุมชัดเจน แต่เมื่อใช้ระยะมีเครื่องหมาย สมการเดียวกันครอบคลุมวัตถุเสมือน ภาพเสมือน และผิวที่มี $R<0$

\subsection{ผิวราบเป็นกรณีพิเศษและเชื่อมกับความลึกปรากฏ}

ผิวราบมีรัศมีความโค้งอนันต์

\[
R\to\infty
\]

ดังนั้น

\[
\frac{n_2-n_1}{R}\to0
\]

สมการผิวหักเหกลายเป็น

\[
\frac{n_1}{p}+\frac{n_2}{q}=0
\]

ย้ายพจน์

\[
\frac{n_2}{q}=-\frac{n_1}{p}
\]

จึงได้

\[
q=-\frac{n_2}{n_1}p
\]

เครื่องหมายลบบอกว่าภาพอยู่ด้านเดียวกับวัตถุและเป็นภาพเสมือน ขนาดความลึกปรากฏคือ

\[
d_{\text{app}}=|q|=\frac{n_2}{n_1}p
\]

เมื่อ $p=d$ จะกลับไปได้สูตรความลึกปรากฏที่อนุมานก่อนหน้า

% =========================================================
\section{เลนส์บาง}

\subsection{เลนส์รวมแสงและเลนส์กระจายแสง}

เลนส์จะรวมแสงหรือกระจายแสงขึ้นกับทั้งรูปร่างและดรรชนีสัมพัทธ์

\[
n_{\text{rel}}\equiv\frac{n}{n_m}
\]

โดย $n$ เป็นดรรชนีของเลนส์และ $n_m$ เป็นดรรชนีของตัวกลางรอบเลนส์ ในอากาศมักมี $n>n_m$ เลนส์นูนทั่วไปจึงรวมแสงและเลนส์เว้าทั่วไปจึงกระจายแสง แต่ถ้านำเลนส์ไปแช่ในตัวกลางที่มีดรรชนีสูงกว่าเลนส์ บทบาทอาจกลับกัน

\subsection{อนุมานสมการช่างทำเลนส์จากผิวหักเหสองผิว}

ให้แสงเดินจากซ้ายไปขวา เลนส์ดรรชนี $n$ อยู่ในตัวกลางดรรชนี $n_m$ ผิวแรกมีรัศมี $R_1$ และผิวที่สองมีรัศมี $R_2$ ตามข้อตกลงคาร์ทีเซียน สำหรับเลนส์นูนสองด้านในอากาศจะมี

\[
R_1>0
\qquad
R_2<0
\]

ที่ผิวแรก แสงเดินจาก $n_m$ ไป $n$ ใช้สมการผิวหักเห

\[
\frac{n_m}{p}
+
\frac{n}{q_1}
=
\frac{n-n_m}{R_1}
\tag{1}
\]

ถ้าไม่มีผิวที่สอง รังสีหลังผิวแรกจะมุ่งไปสร้างภาพที่ระยะ $q_1$

สำหรับเลนส์บาง เราละเลยระยะห่างระหว่างสองผิว จุดภาพจากผิวแรกจึงทำหน้าที่เป็นวัตถุของผิวที่สอง แต่รังสีก่อนถึงผิวที่สองกำลังรวมเข้าหาจุดด้านขวา จุดนั้นเป็นวัตถุเสมือนของผิวที่สอง จึงมี

\[
p_2=-q_1
\]

ที่ผิวที่สอง แสงเดินจาก $n$ ไป $n_m$

\[
\frac{n}{p_2}
+
\frac{n_m}{q}
=
\frac{n_m-n}{R_2}
\]

แทน $p_2=-q_1$

\[
-\frac{n}{q_1}
+
\frac{n_m}{q}
=
\frac{n_m-n}{R_2}
\tag{2}
\]

บวกสมการ (1) และ (2)

\[
\frac{n_m}{p}
+
\frac{n}{q_1}
-
\frac{n}{q_1}
+
\frac{n_m}{q}
=
\frac{n-n_m}{R_1}
+
\frac{n_m-n}{R_2}
\]

เพราะ $n_m-n=-(n-n_m)$

\[
n_m\left(\frac{1}{p}+\frac{1}{q}\right)
=
(n-n_m)
\left(\frac{1}{R_1}-\frac{1}{R_2}\right)
\]

หารด้วย $n_m$

\[
\frac{1}{p}+\frac{1}{q}
=
\left(\frac{n}{n_m}-1\right)
\left(\frac{1}{R_1}-\frac{1}{R_2}\right)
\tag{3}
\]

นิยามความยาวโฟกัสภาพโดยให้วัตถุอยู่ไกลอนันต์

\[
p\to\infty
\qquad
q=f
\]

แทนใน (3)

\[
\frac{1}{f}
=
\left(\frac{n}{n_m}-1\right)
\left(\frac{1}{R_1}-\frac{1}{R_2}\right)
\]

\begin{tcolorbox}[title={สมการช่างทำเลนส์}]
\[
\boxed{
\frac{1}{f}
=
\left(\frac{n}{n_m}-1\right)
\left(\frac{1}{R_1}-\frac{1}{R_2}\right)}
\]

ในอากาศ $n_m\approx1$

\[
\boxed{
\frac{1}{f}
=
(n-1)
\left(\frac{1}{R_1}-\frac{1}{R_2}\right)}
\]
\end{tcolorbox}

\subsection{อนุมานสมการเลนส์บาง}

จากสมการ (3) ด้านขวาขึ้นกับสมบัติของเลนส์และตัวกลางเท่านั้น และจากนิยามโฟกัส ด้านขวาเท่ากับ $1/f$ ดังนั้น

\begin{tcolorbox}[title={สมการเลนส์บาง}]
\[
\boxed{\frac{1}{p}+\frac{1}{q}=\frac{1}{f}}
\]
\end{tcolorbox}

สมการเลนส์จึงไม่ได้ถูกตั้งขึ้นเพราะมีรูปเหมือนสมการกระจก แต่ได้จากการใช้สมการผิวหักเหสองครั้งแล้วให้พจน์ภาพกลางหักล้างกัน

\subsection{อนุมานกำลังขยายของเลนส์บาง}

รังสีที่ผ่านบริเวณศูนย์กลางเชิงแสงของเลนส์บางพบผิวสองด้านที่เกือบขนานกัน การเบี่ยงที่ผิวแรกและผิวที่สองจึงเกือบหักล้างกัน เราประมาณว่ารังสีเส้นนี้เดินเป็นเส้นตรง

จากสามเหลี่ยมด้านวัตถุและภาพซึ่งมีมุมที่ศูนย์กลางเชิงแสงเท่ากัน

\[
\frac{h}{p}=-\frac{h'}{q}
\]

เครื่องหมายลบเกิดจากภาพจริงอยู่คนละด้านของแกนกับยอดวัตถุ จัดรูป

\[
\frac{h'}{h}=-\frac{q}{p}
\]

ใช้คำนิยาม $m=h'/h$

\begin{tcolorbox}[title={กำลังขยายของเลนส์บาง}]
\[
\boxed{m=-\frac{q}{p}}
\]
\end{tcolorbox}

\subsection{รังสีสำคัญของเลนส์มาจากอะไร}

\begin{enumerate}
    \item รังสีขนานแกนเทียบเท่าวัตถุที่ $p\to\infty$ สมการเลนส์ให้ $q=f$ จึงผ่านโฟกัสภาพ
    \item จากการย้อนกลับของรังสี รังสีที่มุ่งผ่านโฟกัสวัตถุจะออกขนานแกน
    \item รังสีผ่านศูนย์กลางเชิงแสงเบี่ยงน้อยมากสำหรับเลนส์บาง จึงประมาณเป็นเส้นตรง
\end{enumerate}

\subsection{กรณีภาพของเลนส์นูนและเลนส์เว้า}

\begin{center}
\begin{tabular}{|c|c|c|c|}
\hline
ตำแหน่งวัตถุ & ตำแหน่งภาพ & ลักษณะภาพ & ขนาด \\
\hline
$p>2f$ & $f<q<2f$ & จริง กลับหัว & เล็กลง \\
\hline
$p=2f$ & $q=2f$ & จริง กลับหัว & เท่าเดิม \\
\hline
$f<p<2f$ & $q>2f$ & จริง กลับหัว & ใหญ่ขึ้น \\
\hline
$p=f$ & $q\to\infty$ & รังสีออกขนาน & ไม่กำหนด \\
\hline
$0<p<f$ & $q<0$ & เสมือน ตั้งตรง & ใหญ่ขึ้น \\
\hline
\end{tabular}
\end{center}

สำหรับเลนส์เว้า $f<0$ และวัตถุจริง $p>0$

\[
\frac{1}{q}=\frac{1}{f}-\frac{1}{p}
\]

ด้านขวาเป็นลบ จึงมี $q<0$ ภาพเป็นภาพเสมือน จาก $m=-q/p$ จึงมี $m>0$ และภาพตั้งตรง

\begin{examplebox}{ตัวอย่าง: เลนส์เว้า}
เลนส์เว้ามี $f=-12\,\mathrm{cm}$ วางวัตถุจริงที่ $p=24\,\mathrm{cm}$

\[
\frac{1}{p}+\frac{1}{q}=\frac{1}{f}
\]

\[
\frac{1}{24}+\frac{1}{q}=-\frac{1}{12}
\]

ย้าย $1/24$

\[
\frac{1}{q}
=-\frac{1}{12}-\frac{1}{24}
=-\frac{3}{24}
=-\frac{1}{8}
\]

\[
q=-8\,\mathrm{cm}
\]

\[
m=-\frac{q}{p}
=-\frac{-8}{24}
=\frac{1}{3}
\]

$q<0$ จึงเป็นภาพเสมือน $m>0$ จึงตั้งตรง และมีขนาดหนึ่งในสามของวัตถุ
\end{examplebox}

\subsection{กำลังเลนส์และการอนุมานสูตรเลนส์ชิดกัน}

กำลังเลนส์เป็น \textbf{นิยาม}

\begin{tcolorbox}[title={นิยาม: กำลังเลนส์}]
\[
\boxed{P\equiv\frac{1}{f}}
\]

เมื่อ $f$ ใช้หน่วยเมตร หน่วยของ $P$ คือไดออปเตอร์

\[
1\,\mathrm{D}=1\,\mathrm{m^{-1}}
\]
\end{tcolorbox}

ต่อไปจะอนุมานการบวกกำลัง ไม่ใช่ประกาศสูตร ให้เลนส์บางสองอันชิดกัน

สำหรับเลนส์ที่หนึ่ง

\[
\frac{1}{p}+\frac{1}{q_1}=\frac{1}{f_1}
\tag{4}
\]

ภาพจากเลนส์แรกเป็นวัตถุของเลนส์ที่สอง และเพราะเลนส์ชิดกัน

\[
p_2=-q_1
\]

สำหรับเลนส์ที่สอง

\[
\frac{1}{p_2}+\frac{1}{q}=\frac{1}{f_2}
\]

\[
-\frac{1}{q_1}+\frac{1}{q}=\frac{1}{f_2}
\tag{5}
\]

บวก (4) และ (5) พจน์ $q_1$ หักล้าง

\[
\frac{1}{p}+\frac{1}{q}
=
\frac{1}{f_1}+\frac{1}{f_2}
\]

ถ้าแทนเลนส์สองอันด้วยเลนส์สมมูล $f_{\text{eq}}$

\[
\frac{1}{p}+\frac{1}{q}=\frac{1}{f_{\text{eq}}}
\]

เปรียบเทียบสองสมการ

\begin{tcolorbox}[title={เลนส์บางชิดกัน}]
\[
\boxed{
\frac{1}{f_{\text{eq}}}
=
\frac{1}{f_1}+\frac{1}{f_2}}
\]

หรือ

\[
\boxed{P_{\text{eq}}=P_1+P_2}
\]
\end{tcolorbox}

สูตรนี้ใช้เมื่อระยะห่างระหว่างเลนส์เล็กจนละเลยได้ ถ้าเลนส์ห่างกันต้องทำภาพต่อภาพ

% =========================================================
\section{ระบบทัศนศาสตร์หลายชิ้น}

\subsection{วิธีทำแบบภาพต่อภาพ}

เมื่อมีหลายอุปกรณ์ ให้เดินตามรังสีจริงทีละชิ้น

\begin{enumerate}
    \item ใช้อุปกรณ์ชิ้นแรกหาภาพที่หนึ่ง
    \item พิจารณาว่ารังสีก่อนถึงชิ้นถัดไปกำลังแยกจากหรือรวมเข้าหาภาพที่หนึ่ง
    \item กำหนด $p$ ของชิ้นถัดไปเป็นบวกสำหรับวัตถุจริง และลบสำหรับวัตถุเสมือน
    \item ทำซ้ำจนถึงอุปกรณ์สุดท้าย
\end{enumerate}

ถ้าความสูงเปลี่ยนจาก $h_0$ เป็น $h_1,h_2,\ldots,h_N$

\[
m_1=\frac{h_1}{h_0}
\]

\[
m_2=\frac{h_2}{h_1}
\]

คูณกัน

\[
m_1m_2
=
\frac{h_1}{h_0}\frac{h_2}{h_1}
=
\frac{h_2}{h_0}
\]

ทำต่อไป พจน์กลางหักล้างทั้งหมด

\begin{tcolorbox}[title={กำลังขยายรวม}]
\[
\boxed{m_{\text{total}}=m_1m_2\cdots m_N}
\]
\end{tcolorbox}

\subsection{ระบบเลนส์สองอันที่แสงเข้าและออกขนาน}

ระบบที่รับรังสีขนานและส่งออกรังสีขนานเรียกว่า \textbf{afocal}

รังสีขนานเข้าหาเลนส์แรกจะมุ่งไปยังโฟกัสภาพของเลนส์แรก ตำแหน่งนี้วัดจากเลนส์แรกเป็นระยะเชิงเครื่องหมาย $f_1$

เพื่อให้รังสีออกจากเลนส์ที่สองขนาน จุดเดียวกันต้องเป็นโฟกัสวัตถุของเลนส์ที่สอง ซึ่งอยู่ก่อนเลนส์ที่สองเป็นระยะ $f_2$

ถ้าเลนส์ห่างกัน $d$ การให้จุดโฟกัสทั้งสองตรงกันเขียนได้เป็น

\[
f_1=d-f_2
\]

จึงได้

\begin{tcolorbox}[title={ระยะห่างของระบบ afocal}]
\[
\boxed{d=f_1+f_2}
\]
\end{tcolorbox}

สำหรับระบบกาลิเลโอ เลนส์แรกเป็นเลนส์เว้า $f_1<0$ และเลนส์ที่สองเป็นเลนส์นูน $f_2>0$

\[
d=f_2-|f_1|
\]

\subsection{อนุมานอัตราส่วนเส้นผ่านศูนย์กลางลำแสง}

ให้รังสีขอบลำแสงอยู่ห่างแกน $D_{\text{in}}/2$ ที่เลนส์แรก และ $D_{\text{out}}/2$ ที่เลนส์ที่สอง สามเหลี่ยมทั้งสองมีมุมของรังสีที่จุดโฟกัสร่วมเท่ากัน

\[
\frac{D_{\text{in}}/2}{|f_1|}
=
\frac{D_{\text{out}}/2}{|f_2|}
\]

ตัด $1/2$

\[
\frac{D_{\text{in}}}{|f_1|}
=
\frac{D_{\text{out}}}{|f_2|}
\]

จัดรูป

\begin{tcolorbox}[title={อัตราส่วนขนาดลำแสงในระบบ afocal}]
\[
\boxed{
\frac{D_{\text{out}}}{D_{\text{in}}}
=
\left|\frac{f_2}{f_1}\right|}
\]
\end{tcolorbox}

\subsection{อนุมานอัตราส่วนความเข้ม}

สำหรับลำแสงที่ความเข้มสม่ำเสมอ กำลังแสง $\mathcal P$ นิยามจาก

\[
\mathcal P=IA
\]

ถ้าละเลยการดูดกลืนและการสะท้อน กำลังแสงคงตัว

\[
I_{\text{in}}A_{\text{in}}
=
I_{\text{out}}A_{\text{out}}
\]

ดังนั้น

\[
\frac{I_{\text{out}}}{I_{\text{in}}}
=
\frac{A_{\text{in}}}{A_{\text{out}}}
\]

พื้นที่หน้าตัดวงกลมคือ

\[
A=\frac{\pi D^2}{4}
\]

จึงได้

\begin{tcolorbox}[title={ความเข้มของลำแสงในระบบอุดมคติ}]
\[
\boxed{
\frac{I_{\text{out}}}{I_{\text{in}}}
=
\left(
\frac{D_{\text{in}}}{D_{\text{out}}}
\right)^2}
\]
\end{tcolorbox}

\subsection{หลักการย้อนกลับของรังสีมาจากอะไร}

กฎการสะท้อน

\[
\theta_i=\theta_r
\]

ไม่เปลี่ยนเมื่อสลับทิศรังสี

กฎของสเนลล์

\[
n_1\sin\theta_1=n_2\sin\theta_2
\]

ก็ยังเป็นสมการเดิมเมื่อสลับตัวกลาง 1 กับ 2 และสลับ $\theta_1$ กับ $\theta_2$ ดังนั้นเส้นทางที่เป็นไปได้จาก $A$ ไป $B$ ย่อมเป็นเส้นทางที่เป็นไปได้จาก $B$ กลับไป $A$ ในระบบตัวกลางนิ่ง

\begin{examplebox}{ตัวอย่าง: ระบบขยายลำแสง}
เลนส์เว้ามี $f_1=-5\,\mathrm{cm}$ และเลนส์นูนมี $f_2=20\,\mathrm{cm}$

ระยะห่างสำหรับระบบ afocal

\[
d=f_1+f_2
=-5+20
=15\,\mathrm{cm}
\]

อัตราส่วนเส้นผ่านศูนย์กลาง

\[
\frac{D_{\text{out}}}{D_{\text{in}}}
=
\left|\frac{20}{-5}\right|
=4
\]

พื้นที่จึงเพิ่ม $4^2=16$ เท่า ถ้ากำลังแสงคงตัว

\[
\frac{I_{\text{out}}}{I_{\text{in}}}
=
\frac{1}{16}
\]
\end{examplebox}

% =========================================================
\section{การเปลี่ยนแปลงเล็กน้อยของความยาวโฟกัส}

\subsection{หลักทั่วไป: เริ่มจากกำลังเลนส์ก่อน}

ให้

\[
P=\frac{1}{f}
\]

หาอนุพันธ์

\[
dP=-\frac{1}{f^2}df
\]

หารด้วย $P=1/f$

\[
\frac{dP}{P}=-\frac{df}{f}
\]

ดังนั้น

\begin{tcolorbox}[title={ความสัมพันธ์ทั่วไปของการเปลี่ยนแปลงเล็กน้อย}]
\[
\boxed{\frac{\delta f}{f}=-\frac{\delta P}{P}}
\]
\end{tcolorbox}

กำลังเพิ่มหนึ่งเปอร์เซ็นต์ทำให้ความยาวโฟกัสลดประมาณหนึ่งเปอร์เซ็นต์

\subsection{เมื่อรัศมีทั้งสองเปลี่ยนทั่วไป}

ในอากาศ

\[
P=(n-1)
\left(
\frac{1}{R_1}-\frac{1}{R_2}
\right)
\]

เมื่อ $n$ คงที่

\[
\delta P
=(n-1)
\left(
-\frac{\delta R_1}{R_1^2}
+
\frac{\delta R_2}{R_2^2}
\right)
\]

ดังนั้น

\[
\frac{\delta f}{f}
=-
\frac{
-\delta R_1/R_1^2+\delta R_2/R_2^2
}{
1/R_1-1/R_2
}
\]

หรือ

\begin{tcolorbox}[title={รัศมีสองผิวเปลี่ยนอย่างอิสระ}]
\[
\boxed{
\frac{\delta f}{f}
=
\frac{
\delta R_1/R_1^2-\delta R_2/R_2^2
}{
1/R_1-1/R_2
}}
\]
\end{tcolorbox}

สูตรง่าย

\[
\frac{\delta f}{f}=\frac{\delta R}{R}
\]

ใช้ได้เมื่อรูปร่างทั้งเลนส์ขยายหรือหดด้วยสัดส่วนเดียวกัน กล่าวคือ $R_1$ และ $R_2$ ถูกคูณด้วยตัวประกอบเดียวกัน ไม่ใช่เมื่อเปลี่ยนเพียงผิวเดียว

ถ้า

\[
R_1'=(1+\varepsilon)R_1
\qquad
R_2'=(1+\varepsilon)R_2
\]

จะมี

\[
P'=\frac{P}{1+\varepsilon}
\]

ดังนั้น

\[
f'=(1+\varepsilon)f
\]

และสำหรับ $|\varepsilon|\ll1$

\begin{tcolorbox}[title={เลนส์เปลี่ยนขนาดแต่คงรูป}]
\[
\boxed{\frac{\delta f}{f}=\varepsilon=\frac{\delta R}{R}}
\]
\end{tcolorbox}

\subsection{เมื่อดรรชนีหักเหเปลี่ยน}

ถ้ารูปร่างและตัวกลางรอบเลนส์คงที่

\[
P=(n_{\text{rel}}-1)K
\]

โดย

\[
n_{\text{rel}}=\frac{n}{n_m}
\]

และ

\[
K=\frac{1}{R_1}-\frac{1}{R_2}
\]

ดังนั้น

\[
\delta P=K\,\delta n_{\text{rel}}
\]

หารด้วย $P=(n_{\text{rel}}-1)K$

\[
\frac{\delta P}{P}
=
\frac{\delta n_{\text{rel}}}{n_{\text{rel}}-1}
\]

ใช้ $\delta f/f=-\delta P/P$

\begin{tcolorbox}[title={ดรรชนีสัมพัทธ์เปลี่ยน}]
\[
\boxed{
\frac{\delta f}{f}
=-
\frac{\delta n_{\text{rel}}}{n_{\text{rel}}-1}}
\]
\end{tcolorbox}

ในอากาศ $n_m\approx1$ จึงมี $n_{\text{rel}}\approx n$

\[
\boxed{
\frac{\delta f}{f}
=-\frac{\delta n}{n-1}}
\]

% =========================================================
\section{ความคลาดของกระจกและเลนส์}

สมการที่อนุมานก่อนหน้าใช้รังสีใกล้แกนและเลนส์บาง ถ้าออกนอกเงื่อนไขจะเกิดความคลาด

\subsection{ความคลาดทรงกลม}

ในการอนุมานสมการกระจกและเลนส์ เราใช้

\[
\sin\theta\approx\tan\theta\approx\theta
\]

การประมาณนี้ดีเมื่อ $\theta$ เล็ก รังสีที่ผ่านขอบอุปกรณ์มีมุมใหญ่กว่า จึงไม่ตัดแกนที่ตำแหน่งเดียวกับรังสีใกล้แกน ภาพจึงไม่คม วิธีลดผลคือจำกัดช่องรับแสงหรือใช้ผิวแอสเฟียริก

\subsection{ความคลาดสี}

สำหรับเลนส์รูปร่างคงที่ในอากาศ

\[
\frac{1}{f}=(n-1)K
\]

จึงมี

\[
f=\frac{1}{(n-1)K}
\]

ในแก้วทั่วไป

\[
n_{\text{blue}}>n_{\text{red}}
\]

ดังนั้น

\[
n_{\text{blue}}-1>n_{\text{red}}-1
\]

ตัวส่วนของ $f_{\text{blue}}$ มากกว่า จึงได้

\[
\boxed{f_{\text{blue}}<f_{\text{red}}}
\]

สีต่างกันจึงรวมที่ตำแหน่งต่างกัน สามารถลดความคลาดสีด้วยการประกอบเลนส์วัสดุต่างชนิดให้การกระจายสีชดเชยกัน

\begin{tcolorbox}[title={ข้อจำกัดของสมการในบทนี้}]
สมการกระจกและเลนส์ให้ผลดีเมื่อ

\begin{enumerate}
    \item รังสีทำมุมเล็กกับแกนมุขสำคัญ
    \item เลนส์บางเมื่อเทียบกับระยะวัตถุและระยะภาพ
    \item ช่องเปิดไม่กว้างจนความคลาดทรงกลมเด่น
    \item ผิวและตัวกลางตรงตามแบบจำลอง
\end{enumerate}
\end{tcolorbox}

% =========================================================
\section{แนวทางแก้โจทย์ทัศนศาสตร์เชิงรังสี}

\begin{enumerate}
    \item วาดแกน ตำแหน่งวัตถุ และลูกศรบอกทิศแสงก่อนเขียนเครื่องหมาย
    \item ระบุว่าสมการที่จะใช้เป็นนิยาม กฎพื้นฐาน หรือผลที่อนุมานแล้ว
    \item เขียน $p$, $f$ และ $R$ พร้อมเครื่องหมายตามข้อตกลง
    \item ใช้สมการกระจกหรือเลนส์หา $q$ แล้วแปลเครื่องหมายของ $q$
    \item ใช้ $m=-q/p$ ซึ่งได้จากสามเหลี่ยมคล้าย ไม่ใช่เดาเครื่องหมายภาพ
    \item ระบบหลายชิ้นต้องหาภาพทีละชิ้น และกำหนดวัตถุจริงหรือวัตถุเสมือนใหม่ทุกครั้ง
    \item ตรวจขีดจำกัด เช่น $p\to\infty$ ต้องได้ $q\to f$
    \item ตรวจหน่วยและตรวจว่า $|\sin\theta|\le1$
\end{enumerate}

\begin{tcolorbox}[title={ข้อผิดพลาดที่พบบ่อย}]
\begin{enumerate}
    \item วัดมุมจากผิวแทนเส้นปกติ
    \item ใช้สูตรมุมเล็กทั้งที่มุมมีค่ามาก
    \item ใส่ $f$ ของเลนส์เว้าหรือกระจกนูนเป็นบวก
    \item ใช้ $q>0$ ทุกกรณีแม้ภาพเป็นภาพเสมือน
    \item ใช้สูตรเลนส์ชิดกันกับเลนส์ที่ห่างกันมาก
    \item ใช้กำลังขยายตามขวางกับวัตถุที่วางตามแกน ทั้งที่ปลายแต่ละจุดมีระยะวัตถุต่างกัน
\end{enumerate}
\end{tcolorbox}

% =========================================================
\section{แบบฝึกหัดท้ายเนื้อหา}

\begin{exercisebox}{แบบฝึกหัดระดับพื้นฐานถึง สอวน.}
\begin{enumerate}
    \item แสงเดินจากอากาศเข้าสู่แก้วดรรชนี $n$ ด้วยมุมตกกระทบ $i$ จงหามุมหักเห
    \item แสงเดินจากแก้วดรรชนี $n$ สู่อากาศ จงหามุมวิกฤต
    \item วัตถุอยู่ใต้น้ำลึก $d$ ผู้สังเกตอยู่ในอากาศและมองเกือบตั้งฉาก จงหาความลึกปรากฏ
    \item กระจกเว้ามีรัศมีความโค้ง $R$ วางวัตถุที่ $p=3R/2$ จงหาตำแหน่งภาพและกำลังขยาย
    \item กระจกนูนมีความยาวโฟกัส $-f$ วางวัตถุที่ระยะ $2f$ จงหาตำแหน่งภาพ
    \item เลนส์นูนมีความยาวโฟกัส $f$ วางวัตถุที่ $p=3f/2$ จงหาตำแหน่งภาพและกำลังขยาย
    \item เลนส์เว้ามีความยาวโฟกัส $-f$ วางวัตถุที่ $p=f$ จงหาตำแหน่งภาพและกำลังขยาย
    \item เลนส์สองอันชิดกันมีความยาวโฟกัส $f_1$ และ $f_2$ จงหาความยาวโฟกัสสมมูล
    \item ระบบ afocal มี $f_1=-a$ และ $f_2=b$ โดย $b>a$ จงหาระยะห่างและอัตราส่วนเส้นผ่านศูนย์กลางลำแสง
    \item เลนส์มี $1/f=(n-1)K$ ถ้า $n$ เพิ่มขึ้นเล็กน้อย $\delta n$ จงหา $\delta f/f$
    \item แสงเข้าสู่ปริซึมหน้าแรกในแนวตั้งฉาก และตกหน้าที่สองด้วยมุม $\beta$ จงหามุมแยกระหว่างสองสีที่มีดรรชนี $n_1$ และ $n_2$
    \item วัตถุเป็นแท่งตามแกนจาก $p=2f$ ถึง $p=3f/2$ หน้าเลนส์นูน จงหาอัตราส่วนความยาวภาพตามแกนต่อความยาววัตถุ
\end{enumerate}
\end{exercisebox}

\subsection*{เฉลยพร้อมที่มาของแต่ละสมการ}

\begin{enumerate}
    \item อากาศมีดรรชนีประมาณ 1 เริ่มจากกฎของสเนลล์
    \[
    1\sin i=n\sin r
    \]
    \[
    \sin r=\frac{\sin i}{n}
    \]
    \[
    \boxed{r=\sin^{-1}\left(\frac{\sin i}{n}\right)}
    \]

    \item ที่มุมวิกฤต มุมหักเหเป็น $90^\circ$
    \[
    n\sin\theta_c=1\sin90^\circ
    \]
    \[
    \sin\theta_c=\frac{1}{n}
    \]
    \[
    \boxed{\theta_c=\sin^{-1}\left(\frac{1}{n}\right)}
    \]

    \item ใช้สูตรความลึกปรากฏโดย $n_1=n$ สำหรับน้ำ และ $n_2\approx1$ สำหรับอากาศ
    \[
    d_{\text{app}}=\frac{n_2}{n_1}d
    =\frac{1}{n}d
    \]
    \[
    \boxed{d_{\text{app}}=\frac{d}{n}}
    \]

    \item จาก $f=R/2$
    \[
    p=\frac{3R}{2}=3f
    \]
    สมการกระจกให้
    \[
    \frac{1}{3f}+\frac{1}{q}=\frac{1}{f}
    \]
    \[
    \frac{1}{q}=\frac{2}{3f}
    \]
    \[
    \boxed{q=\frac{3f}{2}=\frac{3R}{4}}
    \]
    \[
    m=-\frac{q}{p}
    =-\frac{3R/4}{3R/2}
    =-\frac{1}{2}
    \]
    \[
    \boxed{m=-\frac{1}{2}}
    \]

    \item กระจกนูนมี $f_{\text{mirror}}=-f$ และ $p=2f$
    \[
    \frac{1}{2f}+\frac{1}{q}=-\frac{1}{f}
    \]
    \[
    \frac{1}{q}=-\frac{1}{f}-\frac{1}{2f}
    =-\frac{3}{2f}
    \]
    \[
    \boxed{q=-\frac{2f}{3}}
    \]

    \item
    \[
    \frac{1}{3f/2}+\frac{1}{q}=\frac{1}{f}
    \]
    \[
    \frac{2}{3f}+\frac{1}{q}=\frac{1}{f}
    \]
    \[
    \frac{1}{q}=\frac{1}{3f}
    \]
    \[
    \boxed{q=3f}
    \]
    \[
    \boxed{m=-\frac{q}{p}=-\frac{3f}{3f/2}=-2}
    \]

    \item เลนส์เว้ามีความยาวโฟกัส $-f$
    \[
    \frac{1}{f}+\frac{1}{q}=-\frac{1}{f}
    \]
    \[
    \frac{1}{q}=-\frac{2}{f}
    \]
    \[
    \boxed{q=-\frac{f}{2}}
    \]
    \[
    \boxed{m=-\frac{q}{p}=\frac{1}{2}}
    \]

    \item จากการบวกสมการของเลนส์สองอันชิดกัน
    \[
    \boxed{\frac{1}{f_{\text{eq}}}=\frac{1}{f_1}+\frac{1}{f_2}}
    \]

    \item
    \[
    d=f_1+f_2=-a+b
    \]
    \[
    \boxed{d=b-a}
    \]
    และจากสามเหลี่ยมคล้าย
    \[
    \boxed{\frac{D_{\text{out}}}{D_{\text{in}}}
    =\left|\frac{b}{-a}\right|=\frac{b}{a}}
    \]

    \item
    \[
    f=\frac{1}{(n-1)K}
    \]
    ลอการิทึมแล้วหาอนุพันธ์ หรือใช้ผล $\delta f/f=-\delta P/P$
    \[
    \boxed{\frac{\delta f}{f}=-\frac{\delta n}{n-1}}
    \]

    \item รังสีเข้าหน้าแรกตามเส้นปกติจึงไม่เบี่ยง ที่หน้าที่สองใช้กฎของสเนลล์จากปริซึมสู่อากาศ
    \[
    n_j\sin\beta=\sin e_j
    \]
    \[
    e_j=\sin^{-1}(n_j\sin\beta)
    \]
    มุมแยกคือผลต่าง
    \[
    \boxed{
    \Delta e
    =
    \left|
    \sin^{-1}(n_1\sin\beta)
    -
    \sin^{-1}(n_2\sin\beta)
    \right|}
    \]
    สูตรนี้ใช้เมื่อ $n_j\sin\beta\le1$ สำหรับทั้งสองสี

    \item ต้องหาภาพของปลายแต่ละจุดแยกกัน

    ปลายที่ $p_A=2f$
    \[
    \frac{1}{2f}+\frac{1}{q_A}=\frac{1}{f}
    \Rightarrow q_A=2f
    \]

    ปลายที่ $p_B=3f/2$
    \[
    \frac{2}{3f}+\frac{1}{q_B}=\frac{1}{f}
    \Rightarrow q_B=3f
    \]

    ความยาววัตถุ
    \[
    L=2f-\frac{3f}{2}=\frac{f}{2}
    \]

    ความยาวภาพตามแกน
    \[
    L'=3f-2f=f
    \]

    ดังนั้น
    \[
    \boxed{\frac{L'}{L}=\frac{f}{f/2}=2}
    \]
\end{enumerate}

% =========================================================
\section{สรุปสูตรสำคัญและเงื่อนไขการใช้}

\begin{tcolorbox}[title={การสะท้อนและการหักเห}]
\textbf{นิยาม}
\[
n=\frac{c}{v}
\qquad
m=\frac{h'}{h}
\]

\textbf{ผลที่อนุมานจากความถี่คงเดิม}
\[
\lambda=\frac{\lambda_0}{n}
\]

\textbf{กฎจากเวลาเดินทางอยู่กับที่}
\[
\theta_i=\theta_r
\]
\[
n_1\sin\theta_1=n_2\sin\theta_2
\]

\textbf{ผลต่อเนื่องจากกฎของสเนลล์}
\[
d_{\text{app}}=\frac{n_2}{n_1}d
\quad\text{เฉพาะมุมเล็ก}
\]
\[
\sin\theta_c=\frac{n_2}{n_1}
\quad(n_1>n_2)
\]
\[
\tan\theta_B=\frac{n_2}{n_1}
\]
\[
\delta\approx(n-1)A
\quad\text{เฉพาะปริซึมบาง}
\]
\end{tcolorbox}

\begin{tcolorbox}[title={กระจก ผิวหักเห และเลนส์}]
\textbf{รังสีใกล้แกน}
\[
\frac{1}{p}+\frac{1}{q}=\frac{1}{f}
\]
\[
m=-\frac{q}{p}
\]
\[
f_{\text{mirror}}=\frac{R}{2}
\]

\textbf{ผิวหักเหทรงกลม}
\[
\frac{n_1}{p}+\frac{n_2}{q}
=\frac{n_2-n_1}{R}
\]

\textbf{เลนส์บางในตัวกลาง}
\[
\frac{1}{f}
=
\left(\frac{n}{n_m}-1\right)
\left(\frac{1}{R_1}-\frac{1}{R_2}\right)
\]
\end{tcolorbox}

\begin{tcolorbox}[title={ระบบหลายชิ้นและการเปลี่ยนแปลงเล็กน้อย}]
\textbf{เลนส์ชิดกันเท่านั้น}
\[
P_{\text{eq}}=P_1+P_2+\cdots
\]

\textbf{ภาพต่อภาพ}
\[
m_{\text{total}}=m_1m_2\cdots
\]

\textbf{ระบบ afocal}
\[
d=f_1+f_2
\]
\[
\frac{D_{\text{out}}}{D_{\text{in}}}
=
\left|\frac{f_2}{f_1}\right|
\]
\[
\frac{I_{\text{out}}}{I_{\text{in}}}
=
\left(\frac{D_{\text{in}}}{D_{\text{out}}}\right)^2
\quad\text{เมื่อกำลังแสงคงตัว}
\]

\textbf{การเปลี่ยนแปลงเล็กน้อย}
\[
\frac{\delta f}{f}=-\frac{\delta P}{P}
\]
\[
\frac{\delta f}{f}
=-\frac{\delta n_{\text{rel}}}{n_{\text{rel}}-1}
\quad\text{เมื่อรูปร่างคงที่}
\]
\end{tcolorbox}


\section{ข้อสอบจริง สอวน. ที่ใช้ความรู้เรื่องแสง กระจก และเลนส์}



% ---------------------------------------------------------
\subsection{ปี 2560 ข้อ 20: ระบบเลนส์นูนกับกระจกนูน}

\vspace{1em}

\IfFileExists{img/posn60q20.png}{
\begin{center}
    \includegraphics[width=1\linewidth]{img/posn60q20.png}
\end{center}
}{
\begin{tcolorbox}[title={ตำแหน่งภาพที่แนะนำ}]
ให้ตัดภาพข้อสอบปี 2560 ข้อ 20 แล้วบันทึกเป็น

\[
\texttt{img/posn60q20.png}
\]
\end{tcolorbox}
}

\begin{examplebox}{เฉลยละเอียด}
เลนส์กับกระจกห่างกัน $5\,\mathrm{cm}$ กระจกนูนมีรัศมีความโค้ง

\[
R=15\,\mathrm{cm}
\]

ศูนย์กลางความโค้งของกระจกนูนอยู่หลังกระจกเป็นระยะ $15\,\mathrm{cm}$ จึงอยู่ห่างจากเลนส์

\[
5+15=20\,\mathrm{cm}
\]

ถ้าเลนส์สร้างภาพแรกที่ศูนย์กลางความโค้งของกระจก รังสีทุกเส้นจะตกกระจกตามแนวปกติและสะท้อนย้อนแนวเดิม จากหลักการย้อนกลับ รังสีผ่านเลนส์กลับไปสร้างภาพที่ตำแหน่งวัตถุเดิม

ดังนั้นสำหรับการผ่านเลนส์ครั้งแรก

\[
p=30\,\mathrm{cm}
\]

\[
q=20\,\mathrm{cm}
\]

ใช้สมการเลนส์

\[
\frac{1}{f}
=
\frac{1}{30}
+
\frac{1}{20}
\]

\[
\frac{1}{f}=\frac{1}{12}
\]

\[
\boxed{f=12\,\mathrm{cm}}
\]

\[
\boxed{\text{ตอบ C}}
\]
\end{examplebox}

% ---------------------------------------------------------
\subsection{ปี 2560 ตอนที่ 2 ข้อ 10: เลื่อนเลนส์และเลื่อนฉาก}

\vspace{1em}

\IfFileExists{img/posn60p2q10.png}{
\begin{center}
    \includegraphics[width=1\linewidth]{img/posn60p2q10.png}
\end{center}
}{
\begin{tcolorbox}[title={ตำแหน่งภาพที่แนะนำ}]
ให้ตัดภาพข้อสอบปี 2560 ตอนที่ 2 ข้อ 10 แล้วบันทึกเป็น

\[
\texttt{img/posn60p2q10.png}
\]
\end{tcolorbox}
}

\begin{examplebox}{เฉลยละเอียด}
ให้ระยะวัตถุเดิมเป็น $p$

ครั้งแรก ระยะภาพคือ

\[
q_1=45\,\mathrm{cm}
\]

จึงมี

\[
\frac{1}{f}=\frac{1}{p}+\frac{1}{45}
\]

เมื่อเลื่อนเลนส์ออกห่างวัตถุ $6\,\mathrm{cm}$

\[
p_2=p+6
\]

ฉากเดิมอยู่ห่างตำแหน่งเลนส์ใหม่ $45-6=39\,\mathrm{cm}$ และโจทย์เลื่อนฉากเข้าหาเลนส์อีก $3\,\mathrm{cm}$ จึงได้

\[
q_2=36\,\mathrm{cm}
\]

ครั้งที่สอง

\[
\frac{1}{f}=\frac{1}{p+6}+\frac{1}{36}
\]

เทียบสองสมการ

\[
\frac{1}{p}-\frac{1}{p+6}
=
\frac{1}{36}-\frac{1}{45}
\]

\[
\frac{6}{p(p+6)}=\frac{1}{180}
\]

\[
p(p+6)=1080
\]

รากบวกคือ

\[
p=30\,\mathrm{cm}
\]

แทนในครั้งแรก

\[
\frac{1}{f}=\frac{1}{30}+\frac{1}{45}=\frac{1}{18}
\]

\[
\boxed{f=18\,\mathrm{cm}}
\]
\end{examplebox}

% ---------------------------------------------------------
\subsection{ปี 2561 ข้อ 20: วัตถุและฉากอยู่ห่างกันคงที่}

\vspace{1em}

\IfFileExists{img/posn61q20.png}{
\begin{center}
    \includegraphics[width=1\linewidth]{img/posn61q20.png}
\end{center}
}{
\begin{tcolorbox}[title={ตำแหน่งภาพที่แนะนำ}]
ให้ตัดภาพข้อสอบปี 2561 ข้อ 20 แล้วบันทึกเป็น

\[
\texttt{img/posn61q20.png}
\]
\end{tcolorbox}
}

\begin{examplebox}{เฉลยละเอียด}
วัตถุและฉากห่างกัน

\[
p+q=25\,\mathrm{cm}
\]

เลนส์นูนมี

\[
f=6\,\mathrm{cm}
\]

จากสมการเลนส์

\[
\frac{1}{p}+\frac{1}{q}=\frac{1}{6}
\]

คูณด้วย $pq$

\[
p+q=\frac{pq}{6}
\]

ดังนั้น

\[
pq=6(25)=150
\]

$p$ และ $q$ เป็นรากของ

\[
x^2-25x+150=0
\]

\[
(x-10)(x-15)=0
\]

จึงมีตำแหน่งเลนส์สองตำแหน่ง คือ $p=10$ หรือ $15\,\mathrm{cm}$

ต้องการภาพใหญ่กว่าวัตถุ

\[
|m|=\frac{q}{p}>1
\]

จึงเลือก

\[
p=10\,\mathrm{cm}
\qquad
q=15\,\mathrm{cm}
\]

\[
\boxed{p=10\,\mathrm{cm}}
\]

\[
\boxed{\text{ตอบ C}}
\]
\end{examplebox}

% ---------------------------------------------------------
\subsection{ปี 2561 ตอนที่ 2 ข้อ 10: อุปกรณ์ขยายลำแสง}

\vspace{1em}

\IfFileExists{img/posn61p2q10.png}{
\begin{center}
    \includegraphics[width=1\linewidth]{img/posn61p2q10.png}
\end{center}
}{
\begin{tcolorbox}[title={ตำแหน่งภาพที่แนะนำ}]
ให้ตัดภาพข้อสอบปี 2561 ตอนที่ 2 ข้อ 10 แล้วบันทึกเป็น

\[
\texttt{img/posn61p2q10.png}
\]
\end{tcolorbox}
}

\begin{examplebox}{เฉลยละเอียด}
เลนส์เว้าอยู่ก่อน มี

\[
f_1=-4.0\,\mathrm{cm}
\]

เลนส์นูนมี

\[
f_2=12\,\mathrm{cm}
\]

แสงเข้าและออกขนาน ระบบจึงเป็น afocal จุดโฟกัสเสมือนของเลนส์เว้าต้องตรงกับจุดโฟกัสหน้าของเลนส์นูน

ระยะห่างคือ

\[
d=f_1+f_2
\]

\[
d=-4+12
\]

\[
\boxed{d=8.0\,\mathrm{cm}}
\]

การวาดรังสี ให้รังสีขนานผ่านเลนส์เว้าแล้วแยกออก โดยต่อย้อนหลังผ่านจุดโฟกัสของเลนส์เว้า จากนั้นรังสีเหล่านี้เดินเข้าหาเลนส์นูนราวกับมาจากจุดโฟกัสหน้าของเลนส์นูน จึงออกจากเลนส์นูนเป็นรังสีขนาน
\end{examplebox}

% ---------------------------------------------------------
\subsection{ปี 2562 ข้อ 2: มุมเบี่ยงเบนจากเลนส์เว้า}

\vspace{1em}

\IfFileExists{img/posn62q2.png}{
\begin{center}
    \includegraphics[width=1\linewidth]{img/posn62q2.png}
\end{center}
}{
\begin{tcolorbox}[title={ตำแหน่งภาพที่แนะนำ}]
ให้ตัดภาพข้อสอบปี 2562 ข้อ 2 แล้วบันทึกเป็น

\[
\texttt{img/posn62q2.png}
\]
\end{tcolorbox}
}

\begin{examplebox}{เฉลยละเอียด}
เลนส์เว้ามีความยาวโฟกัส

\[
f_{\text{lens}}=-f
\]

รังสีขนานแกนตกที่ความสูง $h$ หลังผ่านเลนส์จะกระจายออก โดยต่อย้อนหลังผ่านจุดโฟกัสด้านหน้า

จากสามเหลี่ยม

\[
\tan\theta=\frac{h}{f}
\]

เพราะ $\theta$ เป็นมุมเล็กในหน่วยเรเดียน

\[
\tan\theta\approx\theta
\]

จึงได้

\[
\boxed{\theta\approx\frac{h}{f}}
\]
\end{examplebox}

% ---------------------------------------------------------
\subsection{ปี 2562 ข้อ 10: กำลังขยายของเลนส์นูนเมื่อภาพเป็นภาพเสมือน}

\vspace{1em}

\IfFileExists{img/posn62q10.png}{
\begin{center}
    \includegraphics[width=1\linewidth]{img/posn62q10.png}
\end{center}
}{
\begin{tcolorbox}[title={ตำแหน่งภาพที่แนะนำ}]
ให้ตัดภาพข้อสอบปี 2562 ข้อ 10 แล้วบันทึกเป็น

\[
\texttt{img/posn62q10.png}
\]
\end{tcolorbox}
}

\begin{examplebox}{เฉลยละเอียด}
ภาพเสมือนอยู่ห่างเลนส์ $17\,\mathrm{cm}$ จึงมี

\[
q=-17\,\mathrm{cm}
\]

ใช้สมการเลนส์

\[
\frac{1}{p}-\frac{1}{17}=\frac{1}{f}
\]

\[
\frac{1}{p}=\frac{1}{f}+\frac{1}{17}
=\frac{17+f}{17f}
\]

ดังนั้น

\[
p=\frac{17f}{17+f}
\]

ขนาดกำลังขยายคือ

\[
|m|=\left|-\frac{q}{p}\right|
=\frac{17}{p}
\]

แทน $p$

\[
|m|
=
\frac{17}{17f/(17+f)}
\]

\[
\boxed{|m|=1+\frac{17}{f}}
\]

ภาพตั้งตรง จึงมี $m>0$
\end{examplebox}

% ---------------------------------------------------------
\subsection{ปี 2563 ข้อ 24: การแยกสีด้วยปริซึม}

\vspace{1em}

\IfFileExists{img/posn63q24.png}{
\begin{center}
    \includegraphics[width=1\linewidth]{img/posn63q24.png}
\end{center}
}{
\begin{tcolorbox}[title={ตำแหน่งภาพที่แนะนำ}]
ให้ตัดภาพข้อสอบศูนย์ สอวน. ปี 2563 ข้อ 24 แล้วบันทึกเป็น

\[
\texttt{img/posn63q24.png}
\]
\end{tcolorbox}
}

\begin{examplebox}{เฉลยละเอียด}
จากรูป ลำแสงเข้าสู่หน้าปริซึมแรกในแนวตั้งฉาก จึงไม่เบี่ยงเบนที่หน้าแรก

ที่หน้าปริซึมที่สอง เส้นปกติตั้งฉากกับผิวเอียง ขณะที่ผิวเอียงทำมุม $\beta$ กับฐาน ดังนั้นมุมตกกระทบภายในคือ

\[
90^\circ-\beta
\]

สำหรับแสงสีน้ำเงิน ดรรชนีเป็น $\mu_1$

\[
\sin e_1
=
\mu_1\sin(90^\circ-\beta)
=
\mu_1\cos\beta
\]

\[
e_1=\sin^{-1}(\mu_1\cos\beta)
\]

สำหรับแสงสีแดง ดรรชนีเป็น $\mu_2$

\[
\sin e_2
=
\mu_2\cos\beta
\]

\[
e_2=\sin^{-1}(\mu_2\cos\beta)
\]

เพราะ $\mu_1>\mu_2$ จึงมี $e_1>e_2$ และมุมระหว่างลำแสงคือ

\[
\boxed{
\theta
=
\sin^{-1}(\mu_1\cos\beta)
-
\sin^{-1}(\mu_2\cos\beta)}
\]

สมการนี้ใช้ได้ตราบใดที่ $\mu_1\cos\beta\le1$ และ $\mu_2\cos\beta\le1$ เพื่อให้ทั้งสองสีหักเหออกจากปริซึมได้
\end{examplebox}

% ---------------------------------------------------------
\subsection{ปี 2563 ข้อ 25: บริเวณสะท้อนกลับหมดในลูกแก้วทรงกลม}

\vspace{1em}

\IfFileExists{img/posn63q25.png}{
\begin{center}
    \includegraphics[width=1\linewidth]{img/posn63q25.png}
\end{center}
}{
\begin{tcolorbox}[title={ตำแหน่งภาพที่แนะนำ}]
ให้ตัดภาพข้อสอบศูนย์ สอวน. ปี 2563 ข้อ 25 แล้วบันทึกเป็น

\[
\texttt{img/posn63q25.png}
\]
\end{tcolorbox}
}

\begin{examplebox}{เฉลยละเอียด}
ให้ $Q$ เป็นจุดขอบเขตที่รังสีจากแหล่งกำเนิด $P$ ตกผิวด้วยมุมวิกฤต

สำหรับแก้วสู่อากาศ

\[
\sin\theta_c=\frac{1}{\mu}
\]

ในสามเหลี่ยม $POQ$

\[
OP=OQ=R
\]

จึงเป็นสามเหลี่ยมหน้าจั่ว และมุมตกกระทบที่ $Q$ เท่ากับมุมฐาน

\[
\angle PQO=\theta_c
\]

ดังนั้นมุมที่ศูนย์กลางจากแนวตรงข้าม $P$ ไปยังจุดขอบเขตเท่ากับ

\[
2\theta_c
\]

เส้น $AB$ เป็นคอร์ดที่เชื่อมจุดขอบเขตสองด้าน จึงมีความยาว

\[
AB=2R\sin(2\theta_c)
\]

ใช้

\[
\sin\theta_c=\frac{1}{\mu}
\]

และ

\[
\cos\theta_c
=
\sqrt{1-\frac{1}{\mu^2}}
=
\frac{\sqrt{\mu^2-1}}{\mu}
\]

จึงได้

\[
AB
=4R\sin\theta_c\cos\theta_c
\]

\[
\boxed{
AB
=
\frac{4R\sqrt{\mu^2-1}}{\mu^2}}
\]
\end{examplebox}

% ---------------------------------------------------------
\subsection{ปี 2566 ข้อ 15: รัศมีความโค้งกับความยาวโฟกัส}

\vspace{1em}

\IfFileExists{img/posn66q15.png}{
\begin{center}
    \includegraphics[width=1\linewidth]{img/posn66q15.png}
\end{center}
}{
\begin{tcolorbox}[title={ตำแหน่งภาพที่แนะนำ}]
ให้ตัดภาพข้อสอบปี 2566 ข้อ 15 แล้วบันทึกเป็น

\[
\texttt{img/posn66q15.png}
\]
\end{tcolorbox}
}

\begin{examplebox}{เฉลยละเอียด}
โจทย์ให้

\[
\frac{1}{f}=(\mu-1)\frac{1}{R}
\]

จึงเขียนได้ว่า

\[
f=\frac{R}{\mu-1}
\]

เมื่อ $\mu$ คงที่

\[
\Delta f=\frac{\Delta R}{\mu-1}
\]

หารด้วย

\[
f=\frac{R}{\mu-1}
\]

จะได้

\[
\boxed{\frac{\Delta f}{f}=\frac{\Delta R}{R}}
\]

\[
\boxed{\text{ตอบข้อ 1}}
\]
\end{examplebox}

% ---------------------------------------------------------
\subsection{ปี 2566 ข้อ 16: ดรรชนีต่ำสุดของปริซึมสะท้อนกลับหลัง}

\vspace{1em}

\IfFileExists{img/posn66q16.png}{
\begin{center}
    \includegraphics[width=1\linewidth]{img/posn66q16.png}
\end{center}
}{
\begin{tcolorbox}[title={ตำแหน่งภาพที่แนะนำ}]
ให้ตัดภาพข้อสอบปี 2566 ข้อ 16 แล้วบันทึกเป็น

\[
\texttt{img/posn66q16.png}
\]
\end{tcolorbox}
}

\begin{examplebox}{เฉลยละเอียด}
รังสีตกผิวเอียงของปริซึมด้วยมุม

\[
\theta_i=45^\circ
\]

ต้องเกิดการสะท้อนกลับหมด จึงต้องมี

\[
\theta_c\le45^\circ
\]

สำหรับปริซึมสู่อากาศ

\[
\sin\theta_c=\frac{1}{\mu}
\]

ค่าดรรชนีต่ำสุดเกิดเมื่อ

\[
\theta_c=45^\circ
\]

ดังนั้น

\[
\frac{1}{\mu_{\min}}=\sin45^\circ=\frac{1}{\sqrt{2}}
\]

\[
\boxed{\mu_{\min}=\sqrt{2}}
\]

\[
\boxed{\text{ตอบข้อ 2}}
\]
\end{examplebox}

% ---------------------------------------------------------
\subsection{ปี 2566 ข้อ 17: กลับด้านรูปร่างเลนส์เมนิสคัส}

\vspace{1em}

\IfFileExists{img/posn66q17.png}{
\begin{center}
    \includegraphics[width=1\linewidth]{img/posn66q17.png}
\end{center}
}{
\begin{tcolorbox}[title={ตำแหน่งภาพที่แนะนำ}]
ให้ตัดภาพข้อสอบปี 2566 ข้อ 17 แล้วบันทึกเป็น

\[
\texttt{img/posn66q17.png}
\]
\end{tcolorbox}
}

\begin{examplebox}{เฉลยละเอียด}
ข้อนี้ต้องระวังว่าเมื่อกลับด้านเลนส์ เราไม่ได้เพียงสลับชื่อผิว แต่เครื่องหมายของรัศมีเชิงคาร์ทีเซียนกลับด้วย

ให้รัศมีเชิงเครื่องหมายของรูป A เป็น $\rho_1$ และ $\rho_2$ ดังนั้น

\[
\frac{1}{f_A}
=
(\mu-1)
\left(
\frac{1}{\rho_1}-\frac{1}{\rho_2}
\right)
\]

เมื่อหมุนเลนส์กลับด้าน ผิวที่สองเดิมกลายเป็นผิวแรกใหม่ และทิศของศูนย์กลางความโค้งกลับเมื่อเทียบกับทิศแสง จึงมี

\[
\rho_1'=-\rho_2
\qquad
\rho_2'=-\rho_1
\]

ใช้สมการช่างทำเลนส์กับรูป B

\[
\frac{1}{f_B}
=
(\mu-1)
\left(
\frac{1}{\rho_1'}-
\frac{1}{\rho_2'}
\right)
\]

แทน $\rho_1'=-\rho_2$ และ $\rho_2'=-\rho_1$

\[
\frac{1}{f_B}
=
(\mu-1)
\left(
-\frac{1}{\rho_2}
+\frac{1}{\rho_1}
\right)
\]

จัดลำดับพจน์

\[
\frac{1}{f_B}
=
(\mu-1)
\left(
\frac{1}{\rho_1}-\frac{1}{\rho_2}
\right)
=
\frac{1}{f_A}
\]

ดังนั้นการกลับด้านเลนส์บางในตัวกลางเดียวกันไม่เปลี่ยนกำลังเลนส์

\[
\boxed{f_B=f_A}
\]

เขียนตามสัญลักษณ์ที่โจทย์ให้

\[
\boxed{
\frac{1}{f_B}
=
(\mu-1)
\left(
\frac{1}{R_1}-\frac{1}{R_2}
\right)}
\]

\[
\boxed{\text{ตอบข้อ 2}}
\]
\end{examplebox}

% ---------------------------------------------------------
\subsection{ปี 2567 ข้อ 12: รังสีผ่านหยดน้ำและสะท้อนภายในหนึ่งครั้ง}

\vspace{1em}

\IfFileExists{img/posn67q12.png}{
\begin{center}
    \includegraphics[width=1\linewidth]{img/posn67q12.png}
\end{center}
}{
\begin{tcolorbox}[title={ตำแหน่งภาพที่แนะนำ}]
ให้ตัดภาพข้อสอบปี 2567 ข้อ 12 แล้วบันทึกเป็น

\[
\texttt{img/posn67q12.png}
\]
\end{tcolorbox}
}

\begin{examplebox}{เฉลยละเอียด}
รังสีหักเหเข้าสู่หยดน้ำจากมุมตกกระทบ $i$ เป็นมุมหักเห $\theta$

การหักเหครั้งแรกทำให้แนวรังสีเปลี่ยนไป

\[
i-\theta
\]

ที่ผิวด้านหลัง รังสีสะท้อนภายใน มุมระหว่างแนวรังสีก่อนและหลังสะท้อนเทียบกับแนวเดินหน้าเปลี่ยนไป

\[
180^\circ-2\theta
\]

การหักเหออกอีกครั้งเพิ่มการเบี่ยงเบน

\[
i-\theta
\]

มุมเบี่ยงเบนรวมแบบมีทิศจากแนวเดินหน้าเป็น

\[
D
=
(i-\theta)
+
(180^\circ-2\theta)
+
(i-\theta)
\]

\[
D=180^\circ+2i-4\theta
\]

แต่ในรูป $\phi$ เป็นมุมแคบระหว่างแนวแสงเข้าและแนวแสงออก จึงมี

\[
\phi=180^\circ-D
\]

ดังนั้น

\[
\boxed{\phi=4\theta-2i}
\]

หรือในหน่วยเรเดียน

\[
\boxed{\phi=4\theta-2i}
\]

โดยใช้ $i$ และ $\theta$ ในหน่วยเรเดียนเหมือนกัน
\end{examplebox}

% ---------------------------------------------------------
\subsection{ปี 2567 ข้อ 14: ดรรชนีหักเหเปลี่ยนเล็กน้อย}

\vspace{1em}

\IfFileExists{img/posn67q14.png}{
\begin{center}
    \includegraphics[width=1\linewidth]{img/posn67q14.png}
\end{center}
}{
\begin{tcolorbox}[title={ตำแหน่งภาพที่แนะนำ}]
ให้ตัดภาพข้อสอบปี 2567 ข้อ 14 แล้วบันทึกเป็น

\[
\texttt{img/posn67q14.png}
\]
\end{tcolorbox}
}

\begin{examplebox}{เฉลยละเอียด}
ให้

\[
K\equiv
\left(
\frac{1}{R_1}+\frac{1}{R_2}
\right)
\]

ซึ่งคงที่ตามรูปโจทย์

\[
\frac{1}{f}=(\mu-1)K
\]

จึงมี

\[
f=\frac{1}{(\mu-1)K}
\]

หาอนุพันธ์สัมพัทธ์

\[
\frac{\delta f}{f}
=-\frac{\delta\mu}{\mu-1}
\]

ดังนั้น

\[
\boxed{
\frac{\delta f}{f}
=-\frac{\delta\mu}{\mu-1}}
\]
\end{examplebox}

% ---------------------------------------------------------
\subsection{ปี 2568 ข้อ 13: ความเข้มในระบบเลนส์นูนกับเลนส์เว้า}

\vspace{1em}

\IfFileExists{img/posn68q13.png}{
\begin{center}
    \includegraphics[width=1\linewidth]{img/posn68q13.png}
\end{center}
}{
\begin{tcolorbox}[title={ตำแหน่งภาพที่แนะนำ}]
ให้ตัดภาพข้อสอบปี 2568 ข้อ 13 แล้วบันทึกเป็น

\[
\texttt{img/posn68q13.png}
\]
\end{tcolorbox}
}

\begin{examplebox}{เฉลยละเอียด}
ระบบรับลำแสงขนานและส่งออกเป็นลำแสงขนาน จึงเป็นระบบ afocal

โจทย์ให้

\[
f_2=-\frac{1}{2}f_1
\]

อัตราส่วนเส้นผ่านศูนย์กลางลำแสงคือ

\[
\frac{D_2}{D_1}
=
\left|\frac{f_2}{f_1}\right|
=
\frac{1}{2}
\]

พื้นที่หน้าตัดจึงเปลี่ยนเป็น

\[
\frac{A_2}{A_1}
=
\left(\frac{D_2}{D_1}\right)^2
=
\frac{1}{4}
\]

เมื่อกำลังแสงคงตัว

\[
IA=\text{constant}
\]

ดังนั้น

\[
\frac{I_2}{I_1}
=
\frac{A_1}{A_2}
=4
\]

\[
\boxed{\frac{I_2}{I_1}=4}
\]
\end{examplebox}

% ---------------------------------------------------------
\subsection{ปี 2568 ข้อ 14: ภาพของวัตถุที่วางตามแกนเลนส์}

\vspace{1em}

\IfFileExists{img/posn68q14.png}{
\begin{center}
    \includegraphics[width=1\linewidth]{img/posn68q14.png}
\end{center}
}{
\begin{tcolorbox}[title={ตำแหน่งภาพที่แนะนำ}]
ให้ตัดภาพข้อสอบปี 2568 ข้อ 14 แล้วบันทึกเป็น

\[
\texttt{img/posn68q14.png}
\]
\end{tcolorbox}
}

\begin{examplebox}{เฉลยละเอียด}
จุด $A$ อยู่ห่างเลนส์

\[
p_A=2f
\]

สมการเลนส์ให้

\[
\frac{1}{2f}+\frac{1}{q_A}=\frac{1}{f}
\]

\[
q_A=2f
\]

วัตถุยาวตามแกน

\[
AB=0.5f
\]

จุด $B$ อยู่ใกล้เลนส์กว่า จึงมี

\[
p_B=2f-0.5f=1.5f
\]

ใช้สมการเลนส์

\[
\frac{1}{1.5f}+\frac{1}{q_B}=\frac{1}{f}
\]

\[
\frac{2}{3f}+\frac{1}{q_B}=\frac{1}{f}
\]

\[
q_B=3f
\]

ความยาวภาพตามแกนคือ

\[
A'B'=q_B-q_A=3f-2f=f
\]

ดังนั้น

\[
\frac{A'B'}{AB}
=
\frac{f}{0.5f}
\]

\[
\boxed{\frac{A'B'}{AB}=2}
\]

ข้อนี้ต้องหาภาพของปลายทั้งสองแยกกัน ไม่สามารถใช้กำลังขยายตามขวาง $m=-q/p$ กับความยาวตามแกนโดยตรง
\end{examplebox}


\fussy

% Additional-exam prompts are intentionally unboxed.
\renewtcolorbox{exercisebox}[1]{
    blanker,
    title=#1,
    fonttitle=\bfseries,
    coltitle=black,
    breakable,
    left=0pt,
    right=0pt,
    top=0pt,
    bottom=0pt
}

% ---------------------------------------------------------
\subsection{ข้อเพิ่มเติม 1: เลนส์เว้าชิดกระจกเงาราบ}

\begin{exercisebox}{โจทย์ S2-01}
\textit{ที่มา: ข้อสอบคัดเลือก ม.5 ศูนย์เตรียม 12 ส.ค. 2547}

เลนส์เว้ามีขนาดความยาวโฟกัส $20\,\mathrm{cm}$ วางชิดกับกระจกเงาราบ ลำแสงขนานตกกระทบจากทางซ้าย จงหาตำแหน่งของภาพสุดท้าย

\begin{center}
    \includegraphics[width=0.78\linewidth,keepaspectratio]{img/posn_archive/15-01.png}
\end{center}
\end{exercisebox}

\begin{examplebox}{คำตอบและแนวคิดสำคัญ}
เลนส์เว้าทำให้รังสีขนานแยกออกเสมือนมาจากจุดโฟกัสซึ่งอยู่ห่างเลนส์ $20\,\mathrm{cm}$ เมื่อสะท้อนที่กระจกเงาราบซึ่งวางชิดเลนส์ จุดดังกล่าวถูกสะท้อนสมมาตรไปอีกด้านของกระจก ดังนั้นภาพสุดท้ายเป็นภาพเสมือนอยู่หลังกระจก
\[
\boxed{20\,\mathrm{cm}}
\]
\end{examplebox}

% ---------------------------------------------------------
\subsection{ข้อเพิ่มเติม 2: วัตถุใต้แผ่นแก้วเมื่อภาพปรากฏที่ผิวล่าง}

\begin{exercisebox}{โจทย์ S2-02}
\textit{ที่มา: ข้อสอบคัดเลือก ม.4 ศูนย์เตรียม 8 ก.ย. 2545}

มองวัตถุ $O$ ผ่านแผ่นแก้วดรรชนีหักเห $\mu$ หนา $t$ ในแนวทำมุม $\theta$ กับเส้นตั้งฉากต่อผิวแก้ว พบว่าภาพ $I$ อยู่ที่ผิวล่างของแผ่นแก้วพอดี จงหาระยะที่วัตถุอยู่ใต้ผิวล่างของแผ่นแก้ว

\begin{center}
    \includegraphics[width=0.78\linewidth,keepaspectratio]{img/posn_archive/15-02.png}
\end{center}
\end{exercisebox}

\begin{examplebox}{คำตอบและแนวคิดสำคัญ}
ให้มุมภายในแผ่นแก้วเป็น
\[
r=\sin^{-1}\!\left(\frac{\sin\theta}{\mu}\right).
\]
จากเรขาคณิตของระยะเยื้องด้านข้างของรังสีในอากาศและในแผ่นแก้ว จะได้ระยะวัตถุใต้ผิวล่าง
\[
\boxed{d=t\left[1-\frac{\tan r}{\tan\theta}\right]}
=\boxed{t\left[1-\frac{\tan\left(\sin^{-1}(\sin\theta/\mu)\right)}{\tan\theta}\right]}.
\]
\end{examplebox}

% ---------------------------------------------------------
\subsection{ข้อเพิ่มเติม 3: มุมระหว่างรังสีเข้าและออกของรุ้งปฐมภูมิ}

\begin{exercisebox}{โจทย์ S2-03}
\textit{ที่มา: ข้อสอบคัดเลือก ม.4 ศูนย์เตรียม 28 ส.ค. 2554}

รังสีแสงตกกระทบหยดน้ำทรงกลมด้วยมุม $\theta$ หักเหเข้าสู่หยดน้ำ สะท้อนภายในหนึ่งครั้ง แล้วหักเหออก จงหามุมแคบ $\phi$ ระหว่างแนวรังสีออกกับแนวรังสีเข้าในรูปของ $\theta$ และดรรชนีหักเห $\mu$ ของน้ำ

\begin{center}
    \includegraphics[width=0.78\linewidth,keepaspectratio]{img/posn_archive/15-03.png}
\end{center}
\end{exercisebox}

\begin{examplebox}{คำตอบและแนวคิดสำคัญ}
จากกฎของสเนลล์ มุมหักเหภายในคือ
\[
r=\sin^{-1}\!\left(\frac{\sin\theta}{\mu}\right).
\]
รวมการเปลี่ยนทิศจากการหักเหสองครั้งและการสะท้อนหนึ่งครั้ง จะได้
\[
\boxed{\phi=4r-2\theta
=4\sin^{-1}\!\left(\frac{\sin\theta}{\mu}\right)-2\theta}.
\]
ข้อนี้เป็นโจทย์รุ่นเก่าที่ใช้โครงสร้างเดียวกับข้อสอบปี 2567 ข้อ 12 ซึ่งมีเฉลยละเอียดอยู่ก่อนหน้าแล้ว
\end{examplebox}

% ---------------------------------------------------------
\subsection{ข้อเพิ่มเติม 4: ทรงกลมโปร่งใสโฟกัสรังสีที่ผิวตรงข้าม}

\begin{exercisebox}{โจทย์ S2-04}
\textit{ที่มา: ข้อสอบคัดเลือก ม.4 ศูนย์เตรียม 2 ก.ย. 2550}

ทรงกลมโปร่งใสรัศมี $R$ รับลำรังสีขนานใกล้แกน และทำให้รังสีไปผ่านจุด $P$ บนผิวด้านตรงข้ามพอดี จงหาดรรชนีหักเหของวัสดุ โดยใช้การประมาณมุมเล็ก $\sin\theta\approx\theta$

\begin{center}
    \includegraphics[width=0.78\linewidth,keepaspectratio]{img/posn_archive/15-04.png}
\end{center}
\end{exercisebox}

\begin{examplebox}{คำตอบและแนวคิดสำคัญ}
สำหรับรังสีใกล้แกน ใช้กฎของสเนลล์ที่ผิวแรกและเรขาคณิตของรังสีที่มุ่งจากผิวแรกไปยังจุด $P$ จะได้ว่าอัตราส่วนมุมตกกระทบต่อมุมหักเหเข้าใกล้ 2 ดังนั้น
\[
\boxed{\mu\approx2}.
\]
\end{examplebox}

% ---------------------------------------------------------
\subsection{ข้อเพิ่มเติม 5: ระบบเลนส์นูนและเลนส์เว้าห่างกัน}

\begin{exercisebox}{โจทย์ S2-05}
\textit{ที่มา: ข้อสอบคัดเลือก ม.4 ศูนย์เตรียม 2 ก.ย. 2550}

เลนส์นูนมี $f_1=12\,\mathrm{cm}$ และเลนส์เว้ามี $f_2=-10\,\mathrm{cm}$ วางบนแกนเดียวกันห่างกัน $10\,\mathrm{cm}$ วัตถุอยู่ซ้ายเลนส์นูน $48\,\mathrm{cm}$ จงหาตำแหน่ง กำลังขยาย และลักษณะของภาพสุดท้าย
\end{exercisebox}

\begin{examplebox}{คำตอบและแนวคิดสำคัญ}
เลนส์แรกให้
\[
\frac1{12}=\frac1{48}+\frac1{q_1}
\quad\Rightarrow\quad q_1=16\,\mathrm{cm}.
\]
ภาพกลางอยู่ขวาเลนส์เว้า $6\,\mathrm{cm}$ จึงเป็นวัตถุเสมือนของเลนส์เว้า $p_2=-6\,\mathrm{cm}$
\[
-\frac1{10}=-\frac1{6}+\frac1{q_2}
\quad\Rightarrow\quad q_2=15\,\mathrm{cm}.
\]
กำลังขยายรวม
\[
m=\left(-\frac{16}{48}\right)\left(-\frac{15}{-6}\right)=-\frac56.
\]
ดังนั้น
\[
\boxed{q_2=15\,\mathrm{cm}\text{ ทางขวาของเลนส์เว้า},\quad m=-\frac56}
\]
เป็นภาพจริง กลับหัว และเล็กลง
\end{examplebox}

% ---------------------------------------------------------
\subsection{ข้อเพิ่มเติม 6: ระบบเลนส์นูน–เลนส์เว้าแบบขยายลำแสง}

\begin{exercisebox}{โจทย์ S2-06}
\textit{ที่มา: ข้อสอบคัดเลือก ม.5 ศูนย์เตรียม 1 ก.ย. 2556}

ให้ $f_1$ และ $f_2$ เป็นขนาดของความยาวโฟกัสของเลนส์นูนและเลนส์เว้าตามลำดับ โดย $f_1>f_2$ จงหา (1) ระยะห่างที่ทำให้ลำแสงออกขนาน (2) อัตราส่วนความเข้มลำแสงออกต่อลำแสงเข้า เมื่อไม่คิดการสูญเสีย และ (3) เมื่อเลื่อนเลนส์เว้ามาประกบชิดเลนส์นูน รังสีขนานจะตัดกันที่ใด

\begin{center}
    \includegraphics[width=0.96\linewidth,keepaspectratio]{img/posn_archive/15-05.png}
\end{center}
\end{exercisebox}

\begin{examplebox}{คำตอบและแนวคิดสำคัญ}
ใช้ $f_{\rm convex}=+f_1$ และ $f_{\rm concave}=-f_2$
\[
\boxed{d=f_1-f_2}.
\]
อัตราส่วนเส้นผ่านศูนย์กลางลำแสงคือ $D_{\rm out}/D_{\rm in}=f_1/f_2$ ตามทิศการจัดในรูป จึงได้
\[
\boxed{\frac{I_{\rm out}}{I_{\rm in}}=\left(\frac{f_2}{f_1}\right)^2}.
\]
เมื่อเลนส์ชิดกัน
\[
\frac1F=\frac1{f_1}-\frac1{f_2}
\quad\Rightarrow\quad
F=-\frac{f_1f_2}{f_1-f_2}.
\]
ดังนั้นรังสีตัดย้อนหลังทางซ้ายของเลนส์คู่เป็นระยะ
\[
\boxed{\frac{f_1f_2}{f_1-f_2}}.
\]
\end{examplebox}

% ---------------------------------------------------------
\subsection{ข้อเพิ่มเติม 7: ปริซึมดรรชนี $\sqrt3$}

\begin{exercisebox}{โจทย์ S2-07}
\textit{ที่มา: ข้อสอบคัดเลือก ม.4 ศูนย์เตรียม 28 ส.ค. 2554}

ปริซึมตามรูปมีดรรชนีหักเห $n=\sqrt3$ รังสีเข้าในแนวตั้งฉากกับผิวแรกและออกจากผิวที่สอง จงหามุม $\theta$ ตามรูป

\begin{center}
    \includegraphics[width=0.78\linewidth,keepaspectratio]{img/posn_archive/15-06.png}
\end{center}
\end{exercisebox}

\begin{examplebox}{คำตอบและแนวคิดสำคัญ}
ที่ผิวแรกไม่มีการหักเห มุมตกกระทบภายในที่ผิวที่สองเท่ากับ $30^\circ$ ใช้กฎของสเนลล์
\[
\sqrt3\sin30^\circ=\sin e=\frac{\sqrt3}{2}
\Rightarrow e=60^\circ.
\]
จากรูป $\theta=e-30^\circ$ จึงได้
\[
\boxed{\theta=30^\circ}.
\]
\end{examplebox}

% ---------------------------------------------------------
\subsection{ข้อเพิ่มเติม 8: เปรียบเทียบกำลังขยายเมื่อวางวัตถุสองตำแหน่ง}

\begin{exercisebox}{โจทย์ S2-08}
\textit{ที่มา: ข้อสอบคัดเลือก ม.5 ศูนย์เตรียม 26 ส.ค. 2555}

เมื่อวางวัตถุห่างเลนส์ $a$ ตามรูป ก. จะได้ภาพเสมือนซึ่งมีขนาดเป็นสองเท่าของภาพจริงที่ได้เมื่อวางวัตถุห่างเลนส์ $b$ ตามรูป ข. จงหาความยาวโฟกัสของเลนส์

\begin{center}
    \includegraphics[width=0.96\linewidth,keepaspectratio]{img/posn_archive/15-07.png}
\end{center}
\end{exercisebox}

\begin{examplebox}{คำตอบและแนวคิดสำคัญ}
ใช้สมการเลนส์และกำลังขยาย $m=-q/p$ แยกสองกรณี โดยคำนึงถึงเครื่องหมายของภาพเสมือนและภาพจริง แล้วใช้เงื่อนไขขนาดภาพกรณีแรกเป็นสองเท่าของกรณีที่สอง จะได้
\[
\boxed{f=2a-b}.
\]
\end{examplebox}

% ---------------------------------------------------------
\subsection{ข้อเพิ่มเติม 9: มุมต่ำสุดของปริซึมเพื่อให้เกิดการสะท้อนกลับหมด}

\begin{exercisebox}{โจทย์ S2-09}
\textit{ที่มา: ข้อสอบคัดเลือก ม.5 ศูนย์เตรียม 26 ส.ค. 2555}

ปริซึมแก้วดรรชนี $n$ ตามรูป ต้องมีมุม $\theta$ อย่างน้อยเท่าใดจึงทำให้แสงไม่ทะลุด้าน $AC$ และเมื่อใช้มุมต่ำสุดนั้น แนวแสงที่ออกทางด้าน $BC$ ทำมุม $\phi$ กับแนว $PQ$ เท่าใด

\begin{center}
    \includegraphics[width=0.78\linewidth,keepaspectratio]{img/posn_archive/15-08.png}
\end{center}
\end{exercisebox}

\begin{examplebox}{คำตอบและแนวคิดสำคัญ}
ค่าต่ำสุดเกิดเมื่อรังสีที่ด้าน $AC$ ตกด้วยมุมวิกฤต
\[
\sin\theta_c=\frac1n.
\]
ใช้ความสัมพันธ์มุมจากรูปหา $\theta_{\min}$ แล้วใช้กฎของสเนลล์ที่ด้าน $BC$ จะได้รูปคำตอบของมุมออก
\[
\boxed{\phi=90^\circ+\sin^{-1}\!\left(n\cos2\theta_{\min}\right)}.
\]
ต้องตรวจด้วยว่าอาร์กิวเมนต์ของ $\sin^{-1}$ ไม่เกิน 1
\end{examplebox}

% ---------------------------------------------------------
\subsection{ข้อเพิ่มเติม 10: กำลังขยายของกระจกเว้าเมื่อได้ภาพเสมือน}

\begin{exercisebox}{โจทย์ S2-10}
\textit{ที่มา: ข้อสอบคัดเลือก ม.5 ศูนย์เตรียม 31 ส.ค. 2557}

กระจกเว้ามีความยาวโฟกัส $+f$ วางวัตถุให้ได้ภาพเสมือนอยู่หลังกระจก $17\,\mathrm{cm}$ จงหากำลังขยายในรูปของ $f$

\begin{center}
    \includegraphics[width=0.78\linewidth,keepaspectratio]{img/posn_archive/15-09.png}
\end{center}
\end{exercisebox}

\begin{examplebox}{คำตอบและแนวคิดสำคัญ}
ให้ $q=-17\,\mathrm{cm}$ จากสมการกระจก
\[
\frac1p-\frac1{17}=\frac1f
\Rightarrow p=\frac{17f}{17+f}.
\]
ดังนั้น
\[
m=-\frac{q}{p}=\frac{17}{p}
=\boxed{1+\frac{17}{f}}.
\]
\end{examplebox}

% ---------------------------------------------------------
\subsection{ข้อเพิ่มเติม 11: การเปลี่ยนแปลงเล็กน้อยของระยะวัตถุและระยะภาพ}

\begin{exercisebox}{โจทย์ S2-11}
\textit{ที่มา: ข้อสอบคัดเลือก ม.5 ศูนย์เตรียม 6 ก.ย. 2546}

จากรูป เมื่อลดระยะวัตถุลงเล็กน้อย ระยะภาพเพิ่มขึ้นเล็กน้อย จงหาความสัมพันธ์ระหว่าง $u$ กับ $v$ ที่ทำให้ $-\delta u=4\delta v$

\begin{center}
    \includegraphics[width=0.78\linewidth,keepaspectratio]{img/posn_archive/15-10.png}
\end{center}
\end{exercisebox}

\begin{examplebox}{คำตอบและแนวคิดสำคัญ}
ภายใต้ข้อตกลงที่ $u,v$ เป็นขนาดระยะบวก สมการเลนส์
\[
\frac1u+\frac1v=\frac1f
\]
ให้
\[
-\frac{\delta u}{u^2}-\frac{\delta v}{v^2}=0.
\]
จึงมี $-\delta u=(u^2/v^2)\delta v$ และเงื่อนไขโจทย์ให้
\[
\boxed{u=2v}.
\]
เฉลยต้นฉบับของแหล่งอ้างอิงพิมพ์ความสัมพันธ์แบบพหุนาม $u^2+8uv+4v^2=0$ ซึ่งใช้ข้อตกลงเครื่องหมายต่างจากบทนี้ จึงควรระบุข้อตกลงก่อนเปรียบเทียบคำตอบ
\end{examplebox}

% ---------------------------------------------------------
\subsection{ข้อเพิ่มเติม 12: ปริซึมหน้าจั่วที่รังสีผ่านอย่างสมมาตร}

\begin{exercisebox}{โจทย์ S2-12}
\textit{ที่มา: ข้อสอบคัดเลือก ม.5 ศูนย์เตรียม 6 ก.ย. 2546}

ปริซึมรูปสามเหลี่ยมหน้าจั่วมีมุมยอด $\alpha$ และมีลำแสงผ่านอย่างสมมาตรดังรูป มุมที่รังสีภายนอกทำกับฐานเป็น $\beta$ จงหาดรรชนีหักเหของแก้วในรูปของ $\alpha$ และ $\beta$

\begin{center}
    \includegraphics[width=0.78\linewidth,keepaspectratio]{img/posn_archive/15-11.png}
\end{center}
\end{exercisebox}

\begin{examplebox}{คำตอบและแนวคิดสำคัญ}
ความสมมาตรให้มุมหักเหภายในแต่ละด้านเท่ากับ $\alpha/2$ ส่วนมุมตกกระทบภายนอกวัดจากเส้นปกติเท่ากับ $\alpha/2+\beta$ ดังนั้น
\[
\boxed{n=\frac{\sin(\alpha/2+\beta)}{\sin(\alpha/2)}}.
\]
\end{examplebox}

% ---------------------------------------------------------
\subsection{ข้อเพิ่มเติม 13: เส้นทางแสงในแท่งแก้วสี่เหลี่ยม}

\begin{exercisebox}{โจทย์ S2-13}
\textit{ที่มา: ข้อสอบคัดเลือก ม.4 ศูนย์เตรียม 6 ก.ย. 2546}

แท่งแก้วสี่เหลี่ยมผืนผ้ายาว $L$ กว้าง $1\,\mathrm{mm}$ มีดรรชนี $1.5$ แสงตกด้าน $AB$ ที่กึ่งกลางด้วยมุมตกกระทบ $\theta=\sin^{-1}(3/4)$ สะท้อนที่ด้าน $AD$ 3 ครั้งและด้าน $BC$ 2 ครั้ง ก่อนออกที่กึ่งกลางด้าน $CD$ จงหาระยะทางรวมที่แสงเดินในแก้ว

\begin{center}
    \includegraphics[width=0.78\linewidth,keepaspectratio]{img/posn_archive/15-12.png}
\end{center}
\end{exercisebox}

\begin{examplebox}{คำตอบและแนวคิดสำคัญ}
จากกฎของสเนลล์
\[
\sin r=\frac{\sin\theta}{1.5}=\frac{1}{2}
\Rightarrow r=30^\circ.
\]
องค์ประกอบตามแนวยาวของเส้นทางรวมเท่ากับ $L$ ดังนั้นถ้ารังสีทำมุม $30^\circ$ กับเส้นปกติของด้าน $AB$ ซึ่งเป็นแนวยาว
\[
\ell\cos30^\circ=L.
\]
จึงได้
\[
\boxed{\ell=\frac{L}{\cos30^\circ}=\frac{2\sqrt3}{3}L}.
\]
\end{examplebox}

% ---------------------------------------------------------
\subsection{ข้อเพิ่มเติม 14: หมุนกระจกเว้ารอบจุดกึ่งกลาง}

\begin{exercisebox}{โจทย์ S2-14}
\textit{ที่มา: ข้อสอบคัดเลือก ม.4 ศูนย์เตรียม 12 ก.ย. 2547}

$P$ เป็นจุดกึ่งกลางกระจกเว้าและ $O$ เป็นศูนย์กลางความโค้ง ถ้าหมุนกระจกไปเป็นมุม $\theta$ รอบจุด $P$ แนวรังสีสะท้อนจะเบนจากเดิมกี่เท่า

\begin{center}
    \includegraphics[width=0.78\linewidth,keepaspectratio]{img/posn_archive/15-13.png}
\end{center}
\end{exercisebox}

\begin{examplebox}{คำตอบและแนวคิดสำคัญ}
เมื่อเส้นปกติที่จุดตกกระทบหมุนไป $\theta$ มุมตกกระทบและมุมสะท้อนเปลี่ยนด้านละ $\theta$ แนวรังสีสะท้อนจึงหมุนรวม
\[
\boxed{2\theta}.
\]
\end{examplebox}

% ---------------------------------------------------------
\subsection{ข้อเพิ่มเติม 15: ความสัมพันธ์มุมและมุมเบี่ยงเบนของปริซึม}

\begin{exercisebox}{โจทย์ S2-15}
\textit{ที่มา: ข้อสอบคัดเลือก ม.4 ศูนย์เตรียม 28 ส.ค. 2548}

แสงในอากาศตกกระทบปริซึมดรรชนี $n$ มุมยอด $\alpha$ ด้วยมุม $\theta_1$ ตามรูป ให้ $\delta$ เป็นมุมเบี่ยงเบน จงหา (1) ความสัมพันธ์ของ $\theta_2,\theta_3,\alpha$ (2) ความสัมพันธ์ของ $\theta_1,\theta_4,\delta,\alpha$ และ (3) ค่า $X$ เมื่อเขียน $\delta=\theta_1-\alpha+\sin^{-1}(X)$

\begin{center}
    \includegraphics[width=0.96\linewidth,keepaspectratio]{img/posn_archive/15-14.png}
\end{center}
\end{exercisebox}

\begin{examplebox}{คำตอบและแนวคิดสำคัญ}
จากเรขาคณิตภายในปริซึม
\[
\boxed{\alpha=\theta_2+\theta_3}.
\]
ตามนิยามมุมในรูป
\[
\boxed{\delta=\frac{\pi}{2}+\theta_1-\theta_4-\alpha}.
\]
ใช้กฎของสเนลล์ที่ผิวแรกหา $\theta_2$ แล้วใช้ $\theta_3=\alpha-\theta_2$ ที่ผิวสอง จะได้
\[
\boxed{X=n\sin\!\left[\alpha-\sin^{-1}\!\left(\frac{\sin\theta_1}{n}\right)\right]}.
\]
\end{examplebox}

% ---------------------------------------------------------
\subsection{ข้อเพิ่มเติม 16: การเลื่อนเลนส์สองตำแหน่งและขนาดภาพ}

\begin{exercisebox}{โจทย์ S2-16}
\textit{ที่มา: ข้อสอบคัดเลือก ม.4 ศูนย์เตรียม 28 ส.ค. 2548}

หลอดไฟเล็กสองดวงห่างกัน $s$ อยู่ห่างฉาก $1.0\,\mathrm{m}$ เมื่อวางเลนส์ที่ตำแหน่งหนึ่ง ภาพคมชัดบนฉากห่างกัน $40\,\mathrm{mm}$ และเมื่อเลื่อนเลนส์ไปอีกตำแหน่งซึ่งให้ภาพชัดอีกครั้ง ภาพห่างกัน $2.5\,\mathrm{mm}$ จงหา $s$ และความยาวโฟกัสของเลนส์
\end{exercisebox}

\begin{examplebox}{คำตอบและแนวคิดสำคัญ}
ตำแหน่งเลนส์คู่กันมีระยะวัตถุและระยะภาพสลับกัน จึงมีกำลังขยายสองค่าซึ่งเป็นส่วนกลับกัน ให้
\[
m_1s=40\,\mathrm{mm},\qquad m_2s=2.5\,\mathrm{mm},\qquad m_1m_2=1.
\]
ดังนั้น
\[
s^2=(40)(2.5)\,\mathrm{mm^2}=100\,\mathrm{mm^2}
\Rightarrow \boxed{s=10\,\mathrm{mm}}.
\]
ได้ $m_1=4$ จึงมี $q=4p$ และ $p+q=1000\,\mathrm{mm}$ ให้ $p=200\,\mathrm{mm}$, $q=800\,\mathrm{mm}$
\[
f=\frac{pq}{p+q}=\boxed{160\,\mathrm{mm}}.
\]
\end{examplebox}

% ---------------------------------------------------------
\subsection{ข้อเพิ่มเติม 17: เลนส์เว้าและเลนส์นูนให้ภาพที่อนันต์}

\begin{exercisebox}{โจทย์ S2-17}
\textit{ที่มา: ข้อสอบคัดเลือก ม.4 ศูนย์เตรียม 8 ก.ย. 2545}

จากรูป วัตถุอยู่ซ้ายเลนส์เว้า $30\,\mathrm{cm}$ โดย $f_1=-20\,\mathrm{cm}$ และเลนส์นูนมี $f_2=+20\,\mathrm{cm}$ ต้องวางเลนส์ทั้งสองห่างกันเท่าใดจึงทำให้ภาพสุดท้ายอยู่ที่อนันต์

\begin{center}
    \includegraphics[width=0.78\linewidth,keepaspectratio]{img/posn_archive/15-15.png}
\end{center}
\end{exercisebox}

\begin{examplebox}{คำตอบและแนวคิดสำคัญ}
เลนส์เว้าให้
\[
-\frac1{20}=\frac1{30}+\frac1{q_1}
\Rightarrow q_1=-12\,\mathrm{cm}.
\]
เพื่อให้เลนส์นูนส่งรังสีออกขนาน วัตถุของเลนส์นูนต้องอยู่ที่โฟกัสหน้า ห่างเลนส์นูน $20\,\mathrm{cm}$ จุดภาพเสมือนจากเลนส์แรกอยู่ซ้ายเลนส์แรก $12\,\mathrm{cm}$ ดังนั้น
\[
d+12=20
\Rightarrow\boxed{d=8\,\mathrm{cm}}.
\]
\end{examplebox}

% ---------------------------------------------------------
\subsection{ข้อเพิ่มเติม 18: เลนส์นูนผิวโค้งฉาบเป็นกระจกเว้า}

\begin{exercisebox}{โจทย์ S2-18}
\textit{ที่มา: ข้อสอบคัดเลือก ม.4 ศูนย์เตรียม 3 ก.ย. 2549}

เลนส์นูนบางแบบผิวหน้าราบและผิวหลังโค้งมีความยาวโฟกัส $f$ จากนั้นฉาบผิวโค้งให้เป็นกระจกเว้า แล้วฉายลำแสงขนานจากด้านราบ จงหาตำแหน่งภาพสุดท้ายทางซ้ายของเลนส์ในรูปของ $f$ และดรรชนี $\mu$

\begin{center}
    \includegraphics[width=0.78\linewidth,keepaspectratio]{img/posn_archive/15-16.png}
\end{center}
\end{exercisebox}

\begin{examplebox}{คำตอบและแนวคิดสำคัญ}
จากสมการช่างทำเลนส์ของเลนส์ผิวราบ–ผิวโค้ง หา $R$ ในรูปของ $f$ และ $\mu$ แล้ววิเคราะห์การหักเหเข้า การสะท้อนที่ผิวโค้ง และการหักเหออกตามลำดับ ผลลัพธ์คือภาพสุดท้ายอยู่ทางซ้ายของเลนส์เป็นระยะ
\[
\boxed{x=\left(\frac{\mu-1}{\mu+1}\right)f}.
\]
\end{examplebox}

% ---------------------------------------------------------
\subsection{ข้อเพิ่มเติม 19: การเลื่อนจุดแสงเมื่อหมุนกระจกเงาราบ}

\begin{exercisebox}{โจทย์ S2-19}
\textit{ที่มา: ข้อสอบคัดเลือก ม.4 ศูนย์เตรียม 31 ส.ค. 2551}

รังสีแสงสะท้อนจากกระจกเงาราบไปตกฉากซึ่งตั้งฉากกับรังสีสะท้อนและอยู่ห่างจุดสะท้อน $5.00\,\mathrm{m}$ เมื่อหมุนกระจกไปมุมเล็ก $\alpha$ จุดแสงบนฉากเลื่อน $0.175\,\mathrm{m}$ จงหา $\alpha$ เป็นองศา
\end{exercisebox}

\begin{examplebox}{คำตอบและแนวคิดสำคัญ}
แนวรังสีสะท้อนหมุน $2\alpha$ ดังนั้น
\[
\tan(2\alpha)\approx2\alpha=\frac{0.175}{5.00}=0.0350.
\]
จึงมี
\[
\alpha=0.0175\,\mathrm{rad}\approx\boxed{1.00^\circ}.
\]
\end{examplebox}

% ---------------------------------------------------------
\subsection{ข้อเพิ่มเติม 20: ดรรชนีหักเหจากความยาวคลื่นในสองตัวกลาง}

\begin{exercisebox}{โจทย์ S1-01}
\textit{ที่มา: แสงเชิงเรขาคณิต ชุดที่ 1 ข้อ 1}

แสงในตัวกลาง $A$ มีดรรชนีหักเห $1.50$ และความยาวคลื่น $500\,\mathrm{nm}$ เมื่อเดินทางในตัวกลาง $B$ มีความยาวคลื่น $450\,\mathrm{nm}$ จงหาดรรชนีหักเหของตัวกลาง $B$
\end{exercisebox}

\begin{examplebox}{คำตอบและแนวคิดสำคัญ}
ความถี่คงเดิมและ $\lambda\propto1/n$
\[
\frac{n_B}{n_A}=\frac{\lambda_A}{\lambda_B}
\Rightarrow n_B=1.50\frac{500}{450}=\boxed{1.67}.
\]
\end{examplebox}

% ---------------------------------------------------------
\subsection{ข้อเพิ่มเติม 21: ตำแหน่งวัตถุที่ทำให้ภาพอยู่ที่อนันต์}

\begin{exercisebox}{โจทย์ S1-02}
\textit{ที่มา: แสงเชิงเรขาคณิต ชุดที่ 1 ข้อ 2}

ต้องวางวัตถุที่ตำแหน่งใดหน้ากระจกเว้าจึงทำให้ไม่มีภาพที่ระยะจำกัด
\end{exercisebox}

\begin{examplebox}{คำตอบและแนวคิดสำคัญ}
เมื่อวัตถุอยู่ที่จุดโฟกัส รังสีสะท้อนออกขนานและภาพอยู่ที่อนันต์
\[
\boxed{p=f}.
\]
\end{examplebox}

% ---------------------------------------------------------
\subsection{ข้อเพิ่มเติม 22: ความลึกปรากฏของวัตถุในน้ำ}

\begin{exercisebox}{โจทย์ S1-03}
\textit{ที่มา: แสงเชิงเรขาคณิต ชุดที่ 1 ข้อ 3}

ชายคนหนึ่งมองวัตถุในน้ำตามแนวดิ่ง เห็นภาพสูงกว่าตำแหน่งจริง $10\,\mathrm{cm}$ ดรรชนีน้ำเท่ากับ $4/3$ จงหาว่าภาพอยู่ลึกจากผิวน้ำเท่าใด
\end{exercisebox}

\begin{examplebox}{คำตอบและแนวคิดสำคัญ}
ให้ความลึกจริงเป็น $d$ และความลึกปรากฏเป็น $d'=d/(4/3)=3d/4$ เงื่อนไข $d-d'=10$ ให้ $d/4=10$ ดังนั้น $d=40$ และ
\[
\boxed{d'=30\,\mathrm{cm}}.
\]
\end{examplebox}

% ---------------------------------------------------------
\subsection{ข้อเพิ่มเติม 23: การสะท้อนจากกระจกเงาราบสองแผ่น}

\begin{exercisebox}{โจทย์ S1-04}
\textit{ที่มา: แสงเชิงเรขาคณิต ชุดที่ 1 ข้อ 4}

รังสีขนานตกกระทบกระจกเงาราบสองแผ่นซึ่งทำมุม $\theta$ ต่อกัน หากรังสีสะท้อนสองแนวทำมุม $40^\circ$ และ $30^\circ$ กับแนวดิ่งเดิมตามรูป จงหา $\theta$

\begin{center}
    \includegraphics[width=0.78\linewidth,keepaspectratio]{img/posn_archive/14-01.png}
\end{center}
\end{exercisebox}

\begin{examplebox}{คำตอบและแนวคิดสำคัญ}
ใช้กฎการสะท้อนแยกแต่ละบานแล้วเชื่อมมุมของเส้นปกติกับมุมระหว่างกระจก จะได้
\[
\boxed{\theta=20^\circ}.
\]
\end{examplebox}

% ---------------------------------------------------------
\subsection{ข้อเพิ่มเติม 24: เปรียบเทียบดรรชนีของของเหลวสองชนิด}

\begin{exercisebox}{โจทย์ S1-05}
\textit{ที่มา: แสงเชิงเรขาคณิต ชุดที่ 1 ข้อ 5}

แสงจากอากาศตกของเหลว $a$ ด้วยมุม $60^\circ$ และหักเห $30^\circ$ เมื่อเปลี่ยนเป็นของเหลว $b$ ใช้มุมตกกระทบ $45^\circ$ และยังหักเห $30^\circ$ จงหา $n_b/n_a$
\end{exercisebox}

\begin{examplebox}{คำตอบและแนวคิดสำคัญ}
จากกฎของสเนลล์
\[
n_a=\frac{\sin60^\circ}{\sin30^\circ}=\sqrt3,
\qquad
n_b=\frac{\sin45^\circ}{\sin30^\circ}=\sqrt2.
\]
ดังนั้น
\[
\boxed{\frac{n_b}{n_a}=\sqrt{\frac23}}.
\]
\end{examplebox}

% ---------------------------------------------------------
\subsection{ข้อเพิ่มเติม 25: จุดโฟกัสเมื่อรังสีเข้าสู่แท่งพลาสติก}

\begin{exercisebox}{โจทย์ S1-06}
\textit{ที่มา: แสงเชิงเรขาคณิต ชุดที่ 1 ข้อ 6}

เลนส์นูนมีความยาวโฟกัสในอากาศ $30\,\mathrm{cm}$ วางบนผิวแท่งพลาสติกดรรชนี $1.5$ ลำแสงขนานตกตั้งฉากตามรูป จุดรวมแสงในพลาสติกอยู่ลึกจากผิวเท่าใด

\begin{center}
    \includegraphics[width=0.58\linewidth,keepaspectratio]{img/posn_archive/14-02.png}
\end{center}
\end{exercisebox}

\begin{examplebox}{คำตอบและแนวคิดสำคัญ}
สำหรับผิวราบและรังสีใกล้แกน ระยะภาพเปลี่ยนตามอัตราส่วนดรรชนี
\[
q_{\rm plastic}=\frac{n_{\rm plastic}}{n_{\rm air}}q_{\rm air}
=1.5(30)=\boxed{45\,\mathrm{cm}}.
\]
\end{examplebox}

% ---------------------------------------------------------
\subsection{ข้อเพิ่มเติม 26: ความยาวคลื่นในของเหลวจากมุมวิกฤต}

\begin{exercisebox}{โจทย์ S1-07}
\textit{ที่มา: แสงเชิงเรขาคณิต ชุดที่ 1 ข้อ 7}

มุมวิกฤตของแสงจากของเหลวสู่อากาศเท่ากับ $60^\circ$ ความยาวคลื่นในของเหลวเป็นกี่เท่าของความยาวคลื่นในอากาศ
\end{exercisebox}

\begin{examplebox}{คำตอบและแนวคิดสำคัญ}
มี $\sin\theta_c=1/n$ และ $\lambda_{\rm liquid}/\lambda_{\rm air}=1/n$ ดังนั้น
\[
\boxed{\frac{\lambda_{\rm liquid}}{\lambda_{\rm air}}=\sin60^\circ=\frac{\sqrt3}{2}}.
\]
\end{examplebox}

% ---------------------------------------------------------
\subsection{ข้อเพิ่มเติม 27: หาดรรชนีของปริซึมสามเหลี่ยมมุมฉาก}

\begin{exercisebox}{โจทย์ S1-08}
\textit{ที่มา: แสงเชิงเรขาคณิต ชุดที่ 1 ข้อ 8}

แสงสีเขียวตกตั้งฉากกับด้านหนึ่งของปริซึมสามเหลี่ยมมุมฉากตามรูป และรังสีออกเบนจากแนวเดิม $30^\circ$ จงหาดรรชนีหักเหของปริซึม

\begin{center}
    \includegraphics[width=0.78\linewidth,keepaspectratio]{img/posn_archive/14-03.png}
\end{center}
\end{exercisebox}

\begin{examplebox}{คำตอบและแนวคิดสำคัญ}
ที่ผิวแรกไม่เบี่ยงเบน ใช้เรขาคณิตของปริซึมหามุมตกกระทบที่ผิวเอียง แล้วใช้กฎของสเนลล์ระหว่างปริซึมกับอากาศ จะได้ประมาณ
\[
\boxed{n\approx1.7}.
\]
\end{examplebox}

% ---------------------------------------------------------
\subsection{ข้อเพิ่มเติม 28: สีบนฟองสบู่}

\begin{exercisebox}{โจทย์ S1-09}
\textit{ที่มา: แสงเชิงเรขาคณิต ชุดที่ 1 ข้อ 9}

ริ้วสีต่าง ๆ ที่เห็นบนฟองสบู่เกิดจากปรากฏการณ์ใดเป็นสำคัญ
\end{exercisebox}

\begin{examplebox}{คำตอบและแนวคิดสำคัญ}
เกิดจากการซ้อนทับของแสงที่สะท้อนจากผิวหน้าและผิวหลังของฟิล์มบาง
\[
\boxed{\text{การแทรกสอดของแสง}}
\]
ข้อนี้เป็นทัศนศาสตร์เชิงคลื่น จึงใส่ไว้เป็นโจทย์เสริมท้ายชุด ไม่ควรใช้แบบจำลองรังสีเพียงอย่างเดียว
\end{examplebox}

% ---------------------------------------------------------
\subsection{ข้อเพิ่มเติม 29: เงื่อนไขมุมในเส้นใยนำแสง}

\begin{exercisebox}{โจทย์ S1-10}
\textit{ที่มา: แสงเชิงเรขาคณิต ชุดที่ 1 ข้อ 10}

เส้นใยนำแสงมีแกนดรรชนี $n_1$ และเปลือกดรรชนี $n_2<n_1$ มุมตกกระทบภายในที่รอยต่อเป็น $\theta$ จงเขียนเงื่อนไขที่ทำให้แสงถูกกักอยู่ในแกน

\begin{center}
    \includegraphics[width=0.78\linewidth,keepaspectratio]{img/posn_archive/14-04.png}
\end{center}
\end{exercisebox}

\begin{examplebox}{คำตอบและแนวคิดสำคัญ}
ต้องเกิดการสะท้อนกลับหมด
\[
\boxed{\sin^{-1}\!\left(\frac{n_2}{n_1}\right)<\theta<90^\circ}.
\]
\end{examplebox}

% ---------------------------------------------------------
\subsection{ข้อเพิ่มเติม 30: ระบบเลนส์สามอันรับแสงอาทิตย์}

\begin{exercisebox}{โจทย์ S1-11}
\textit{ที่มา: แสงเชิงเรขาคณิต ชุดที่ 1 ข้อ 11}

เลนส์เว้า $f=-10\,\mathrm{cm}$ สองอันและเลนส์นูน $f=+20\,\mathrm{cm}$ หนึ่งอันวางรับแสงอาทิตย์ตามรูป โดยระยะจากเลนส์เว้าตัวแรกถึงเลนส์นูนเท่ากับ $a$ และจากเลนส์นูนถึงเลนส์เว้าตัวที่สองเท่ากับ $a$ ถ้าลำแสงสุดท้ายเป็นลำขนาน จงหา $a$

\begin{center}
    \includegraphics[width=0.96\linewidth,keepaspectratio]{img/posn_archive/14-05.png}
\end{center}
\end{exercisebox}

\begin{examplebox}{คำตอบและแนวคิดสำคัญ}
ไล่ภาพทีละเลนส์โดยให้วัตถุเริ่มต้นอยู่ที่อนันต์ และกำหนดให้ภาพสุดท้ายอยู่ที่อนันต์ เงื่อนไขภาพต่อภาพให้
\[
\boxed{a=30\,\mathrm{cm}}.
\]
\end{examplebox}

% ---------------------------------------------------------
\subsection{ข้อเพิ่มเติม 31: เลนส์เว้ากับกระจกเว้าให้ภาพกลับที่วัตถุ}

\begin{exercisebox}{โจทย์ S1-12}
\textit{ที่มา: แสงเชิงเรขาคณิต ชุดที่ 1 ข้อ 12}

วัตถุอยู่ซ้ายเลนส์เว้าระยะ $10\,\mathrm{cm}$ เลนส์เว้ามี $f=-10\,\mathrm{cm}$ และกระจกเว้าอยู่ขวาเลนส์ $15\,\mathrm{cm}$ ต้องเลือกความยาวโฟกัสกระจกเท่าใดจึงทำให้ภาพสุดท้ายกลับมาอยู่ที่วัตถุ

\begin{center}
    \includegraphics[width=0.96\linewidth,keepaspectratio]{img/posn_archive/14-06.png}
\end{center}
\end{exercisebox}

\begin{examplebox}{คำตอบและแนวคิดสำคัญ}
ผ่านเลนส์ครั้งแรก
\[
-\frac1{10}=\frac1{10}+\frac1{q_1}
\Rightarrow q_1=-5\,\mathrm{cm}.
\]
รังสีที่ถึงกระจกจึงแยกออกเสมือนมาจากจุดซ้ายเลนส์ $5\,\mathrm{cm}$ ซึ่งอยู่หน้ากระจก $20\,\mathrm{cm}$ เพื่อให้รังสีย้อนแนวเดิม จุดนี้ต้องอยู่ที่ศูนย์กลางความโค้งของกระจก จึงมี $R=20\,\mathrm{cm}$ และ
\[
\boxed{f_{\rm mirror}=10\,\mathrm{cm}}.
\]
ถ้าใช้ข้อตกลงคาร์ทีเซียนที่กระจกเว้ามี $f<0$ จะเขียนเป็น $-10\,\mathrm{cm}$ ซึ่งตรงกับเฉลยต้นฉบับ
\end{examplebox}

% ---------------------------------------------------------
\subsection{ข้อเพิ่มเติม 32: สีที่มีมุมหักเหน้อยที่สุดในปริซึม}

\begin{exercisebox}{โจทย์ S1-13}
\textit{ที่มา: แสงเชิงเรขาคณิต ชุดที่ 1 ข้อ 13}

เมื่อแสงขาวผ่านปริซึม สีใดมีมุมหักเหภายในวัดจากเส้นปกติน้อยที่สุด
\end{exercisebox}

\begin{examplebox}{คำตอบและแนวคิดสำคัญ}
สีม่วงมีดรรชนีหักเหมากที่สุด จึงหักเหเข้าใกล้เส้นปกติมากที่สุดและมีมุมหักเหภายในน้อยที่สุด
\[
\boxed{\text{สีม่วง}}
\]
อย่าสับสนกับมุมเบี่ยงเบน ซึ่งสีม่วงมีค่ามากที่สุด
\end{examplebox}

% ---------------------------------------------------------
\subsection{ข้อเพิ่มเติม 33: ระยะวัตถุของแว่นขยาย}

\begin{exercisebox}{โจทย์ S1-14}
\textit{ที่มา: แสงเชิงเรขาคณิต ชุดที่ 1 ข้อ 14}

แว่นขยายมีความยาวโฟกัส $4\,\mathrm{cm}$ จงหาระยะที่ควรวางตัวหนังสือจากเลนส์ตามคำตอบของชุดข้อสอบต้นฉบับ
\end{exercisebox}

\begin{examplebox}{คำตอบและแนวคิดสำคัญ}
แหล่งอ้างอิงระบุคำตอบ
\[
\boxed{4.76\,\mathrm{cm}}.
\]
อย่างไรก็ตาม สำหรับการใช้แว่นขยายมาตรฐานเพื่อให้ได้ภาพเสมือนที่ระยะใกล้ชัด $D=25\,\mathrm{cm}$ ต้องใช้ $q=-D$ ซึ่งให้
\[
p=\left(\frac1f+\frac1D\right)^{-1}=\frac{100}{29}\approx3.45\,\mathrm{cm}.
\]
ดังนั้นโจทย์ต้นฉบับน่าจะมีเงื่อนไขหรือข้อตกลงเพิ่มเติมที่ไม่ได้ปรากฏในข้อความบนเว็บไซต์ ควรตรวจต้นฉบับเต็มก่อนใช้เป็นข้อสอบให้คะแนน
\end{examplebox}

% ---------------------------------------------------------
\subsection{ข้อเพิ่มเติม 34: สีของวัตถุจากการมองผ่านแผ่นกรอง}

\begin{exercisebox}{โจทย์ S1-15}
\textit{ที่มา: แสงเชิงเรขาคณิต ชุดที่ 1 ข้อ 15}

ดอกไม้ส่องด้วยแสงขาว เห็นเป็นสีแดงเมื่อมองผ่านแผ่นกรองแดง เห็นเป็นสีเขียวผ่านแผ่นกรองเขียว และเห็นเป็นสีเหลืองผ่านแผ่นกรองเหลือง ดอกไม้มีสีใด
\end{exercisebox}

\begin{examplebox}{คำตอบและแนวคิดสำคัญ}
วัตถุต้องสะท้อนทั้งองค์ประกอบสีแดงและสีเขียว ซึ่งรวมกันเป็นสีเหลือง
\[
\boxed{\text{สีเหลือง}}.
\]
ข้อนี้เป็นเรื่องการรับรู้สีและแผ่นกรอง จึงจัดเป็นโจทย์เสริมเกี่ยวกับแสง
\end{examplebox}

\fussy
